%\documentclass[12pt,a4paper]{book}
 \documentclass[a4paper,french]{book}
\usepackage[17pt]{extsizes}

\usepackage[table,dvipsnames,svgnames]{xcolor}
\usepackage{listingsutf8}
%\usepackage[utf8]{inputenc}% gestion des accents (source)
\usepackage[T1]{fontenc}% gestion des accents (PDF)
\usepackage{fontspec}
\usepackage[french]{babel}% gestion du français
\usepackage{textcomp}% caractères additionnels
\usepackage{mathtools,amssymb,amsthm}% packages de l'AMS + mathtools
\usepackage{amsmath}
\usepackage{yhmath}
\usepackage{mathpazo}
\usepackage[scaled=0.95]{helvet}
\usepackage{courier}
%\usepackage{fourier}
%\usepackage[scaled=0.95]{helvet}
%\usepackage{courier}
\usepackage{bm}
\usepackage{gensymb}
\usepackage{graphicx}

\usepackage{enumitem}
\usepackage{pifont}

\usepackage{pstricks}
\usepackage{pstricks-add}
\usepackage{pst-eucl}
\usepackage{pgfplots}
\let\clipbox\relax
%\usepackage{pst-plot,pst-text,pst-tree,pstricks-add}
%\usepackage{pst-node} % required package
\usepackage{pst-all}
\usepackage[most,skins,breakable]{tcolorbox}

\usepackage{euscript}





% ############### Nouvelles commandes ###############
% Coloration des labels des points dans pst-eucl
\makeatletter
\define@key[psset]{pst-eucl}{LabelColor}{\def\psk@LabelColor{#1}}
\psset{LabelColor=black}
\def\Pst@WhichLabel#1{{\color{\psk@LabelColor}%
		\ifx\psk@PointName\@default#1\else\psk@PointName\fi}}
\makeatother

% renomme le symbole #2 en #1#2 et invalide le symbole initial #2:
\def\renamesymbol#1#2{
	\expandafter\let\expandafter\newsym\expandafter=\csname#2\endcsname
	\expandafter\global\expandafter\let\csname#1#2\endcsname=\newsym
	\expandafter\global\expandafter\let\csname#2\endcsname=\origsym
}



\usepackage{ProfCollege}
\renamesymbol{Casio}{Calculatrice}
\usepackage[nonshellescape]{ProfLycee}



%\usepackage[left=1.5cm,right=1.5cm,top=2.5cm,bottom=2.5cm]{geometry}
\usepackage[left=13mm,right=13mm,top=18mm,bottom=18mm,footskip=30pt]{geometry}
\setlength{\headheight}{15pt}

\usepackage{titlesec}
\usepackage{pas-math}
\usepackage{longtable}
\usepackage{tikz,tkz-tab}
\usepackage{amsmath}
\usepackage[right]{eurosym}
\usepackage{sectsty}
\usepackage{tipfr}
\usepackage{setspace}

%pour faire les touches de la calculatrices
\usepackage{menukeys}
\renewmenumacro{\menu}{roundedmenus} % default: menus
\renewmenumacro{\directory}{pathswithfolder} % default: paths
\renewmenumacro{\keys}{shadowedroundedkeys} % default: roundedkeys


\usepackage{multicol}
\usepackage{multirow}
\usepackage{systeme}

%% Format des titres de chapitres, paragraphes, ...
\titleformat{\chapter}[frame]
{\color{red!85!white}\normalsize}%
{\filright\sffamily\Large%
	\enspace Chapitre \thechapter\enspace}%
{8pt}
{\color{red!85!white}\rule{0pt}{30pt}\sffamily\Huge\bfseries\filcenter}
%%%%%%%%%%%%%%%%%%%%
\titlespacing{\chapter}{0pc}{0pc}{2pc}
\titlespacing{\subsection}{2pc}{}{}

\titleformat{\subsection}[block]{\hspace{1em}}{\thesubsection}{1em .}{}
\titleformat{\subsubsection}[block]{\hspace{2em}}{\thesubsection}{1em )}{}

\chapterfont{\color{red}{}}
\sectionfont{\color{green!55!black}{}}
\subsectionfont{\color{cyan!65!black}{}}
\subsubsectionfont{\color{violet}{}}
\renewcommand{\thesection}{\color{green!55!black}\Roman{section}.}
\renewcommand{\thesubsection}{\color{cyan!65!black}\arabic{subsection})}
\renewcommand{\thesubsubsection}{\color{violet}\alph{subsubsection})}
\setcounter{secnumdepth}{3}

\titleformat{\subsection}[block]{\hspace{1em}}{\thesubsection}{1em .}{}
\titleformat{\subsubsection}[block]{\hspace{2em}}{\thesubsection}{1em )}{}

\usepackage[column=A]{cellspace}
\setlength{\cellspacebottomlimit}{4pt}
\setlength{\cellspacetoplimit}{4pt}
\setlength{\parindent}{2pt}
\setlength{\tabcolsep}{0.5cm}


%entête et pieds de pages
\usepackage{fancyhdr}
\pagestyle{fancy}
\fancyhf{}
\usepackage{lastpage}
\renewcommand{\chaptermark}[1]{\markboth {Chap. \thechapter~ - ~#1}{}}
\renewcommand\headrulewidth{1pt}
\fancyhead[LO,RE]{1ère G - spécialité maths}
\fancyhead[RO,LE]{\leftmark}
\renewcommand\footrulewidth{1pt}
\fancyfoot[LO,RE]{Année 2025-2026}
\fancyfoot[RO,LE]{\textbf{Page \thepage/\pageref{LastPage}}}

%Mes boites pour les définitions, propriétés, remarques, ...

%Mes boites pour les définitions, propriétés, remarques, ...

\newcommand\titrebox[1]{\textbf{#1}}
\newtcolorbox{definition}{breakable,colback=red!10!white,
	colframe=red!85!white, fonttitle=\bfseries,enhanced,
	attach boxed title to top left={xshift=5mm, yshift=-2mm},boxed title style={size=small,colback=red!10!white,colframe=red!85!white, arc=0.5mm},
	title=\textcolor{red!85!white}{Définition} , arc=3mm}

\newtcolorbox{definitionTitre}[1]{breakable,colback=red!10!white,
	colframe=red!85!white, fonttitle=\bfseries,enhanced,
	attach boxed title to top left={xshift=5mm, yshift=-2mm},boxed title style={size=small,colback=red!10!white,colframe=red!85!white, arc=0.5mm},
	title=\textcolor{red!85!white}{Définition - #1} , arc=3mm}

\newtcolorbox{definitionPerso}[1]{breakable,colback=red!10!white,
	colframe=red!85!white, fonttitle=\bfseries,enhanced,
	attach boxed title to top left={xshift=5mm, yshift=-2mm},boxed title style={size=small,colback=red!10!white,colframe=red!85!white, arc=0.5mm},
	title=\textcolor{red!85!white}{#1} , arc=3mm}

\newtcolorbox{regle}{breakable,colback=red!10!white,
	colframe=red!85!white, fonttitle=\bfseries,enhanced,
	attach boxed title to top left={xshift=5mm, yshift=-2mm},boxed title style={size=small,colback=red!10!white,colframe=red!85!white, arc=0.5mm},
	title=\textcolor{red!85!white}{Règle} , arc=3mm}

\newtcolorbox{regles}{breakable,colback=red!10!white,
	colframe=red!85!white, fonttitle=\bfseries,enhanced,
	attach boxed title to top left={xshift=5mm, yshift=-2mm},boxed title style={size=small,colback=red!10!white,colframe=red!85!white, arc=0.5mm},
	title=\textcolor{red!85!white}{Règles} , arc=3mm}

\newtcolorbox{definitions}{breakable,colback=red!10!white,
	colframe=red!85!white, fonttitle=\bfseries,enhanced,
	attach boxed title to top left={xshift=5mm, yshift=-2mm},boxed title style={size=small,colback=red!10!white,colframe=red!85!white, arc=0.5mm},
	title=\textcolor{red!85!white}{Définitions} , arc=3mm}

\newtcolorbox{propriete}{breakable,colback=red!10!white,
	colframe=red!85!white, fonttitle=\bfseries,enhanced,
	attach boxed title to top left={xshift=5mm, yshift=-2mm},boxed title style={size=small,colback=red!10!white,colframe=red!85!white, arc=0.5mm},
	title=\textcolor{red!85!white}{Propriété} , arc=3mm}

\newtcolorbox{proprieteTitre}[1]{breakable,colback=red!10!white,
	colframe=red!85!white, fonttitle=\bfseries,enhanced,
	attach boxed title to top left={xshift=5mm, yshift=-2mm},boxed title style={size=small,colback=red!10!white,colframe=red!85!white, arc=0.5mm},
	title=\textcolor{red!85!white}{Propriété - #1} , arc=3mm}

\newtcolorbox{proprietes}{breakable,colback=red!10!white,
	colframe=red!85!white, fonttitle=\bfseries,enhanced,
	attach boxed title to top left={xshift=5mm, yshift=-2mm},boxed title style={size=small,colback=red!10!white,colframe=red!85!white, arc=0.5mm},
	title=\textcolor{red!85!white}{Propriétés} , arc=3mm}

\newtcolorbox{theoreme}{breakable,colback=red!10!white,
	colframe=red!85!white, fonttitle=\bfseries,enhanced,
	attach boxed title to top left={xshift=5mm, yshift=-2mm},boxed title style={size=small,colback=red!10!white,colframe=red!85!white, arc=0.5mm},
	title=\textcolor{red!85!white}{Théorème} , arc=3mm}

\newtcolorbox{theoremeTitre}[1]{breakable,colback=red!10!white,
	colframe=red!85!white, fonttitle=\bfseries,enhanced,
	attach boxed title to top left={xshift=5mm, yshift=-2mm},boxed title style={size=small,colback=red!10!white,colframe=red!85!white, arc=0.5mm},
	title=\textcolor{red!85!white}{Théorème - #1} , arc=3mm}

\newtcolorbox{theoremes}{breakable,colback=red!10!white,
	colframe=red!85!white, fonttitle=\bfseries,enhanced,
	attach boxed title to top left={xshift=5mm, yshift=-2mm},boxed title style={size=small,colback=red!10!white,colframe=red!85!white, arc=0.5mm},
	title=\textcolor{red!85!white}{Théorèmes} , arc=3mm}

\newtcolorbox{bexemple}{
	enhanced,
	breakable,
	fonttitle=\bfseries,
	frame empty,
	nobeforeafter,
	borderline west={0.5mm}{0mm}{cyan!75!black},
	title=\textcolor{cyan!75!black}{\faHubspot\:\:Exemple},
	%title after break=Exemple (suite),
	colback=white,
	extras={enhanced, frame empty},
	detach title,before upper={\tcbtitle\ :\quad },
}

\newtcolorbox{bexemples}{
	enhanced,
	breakable,
	fonttitle=\bfseries,
	frame empty,
	nobeforeafter,
	borderline west={0.5mm}{0mm}{cyan!75!black},
	title=\textcolor{cyan!75!black}{\faHubspot\:\:Exemples},
	%title after break=Exemples (suite),
	colback=white,
	extras={enhanced, frame empty},
	detach title,before upper={\tcbtitle\ :\quad },
}

\newtcolorbox{bexemplet}[1]{
	enhanced,
	breakable,
	fonttitle=\bfseries,
	frame empty,
	borderline west={0.5mm}{0mm}{cyan!75!black},
	title=\textcolor{cyan!75!black}{\faHubspot\:\:Exemple - #1},
	colback=white,
	extras={enhanced, frame empty},
	detach title,before upper={\tcbtitle\ :\quad },
}

\newtcolorbox{brmq}{
	enhanced,
	breakable,
	fonttitle=\bfseries,
	frame empty,
	nobeforeafter,
	borderline west={0.5mm}{0mm}{orange!95!white},
	title=\textcolor{orange!85!white}{\faHandPointRight[regular]\:\:Remarque} ,
	colback=white,
	extras={enhanced, frame empty},
	detach title,before upper={\tcbtitle\ \quad },
}

\newtcolorbox{brmqs}{
	enhanced,
	breakable,
	fonttitle=\bfseries,
	frame empty,
	nobeforeafter,
	borderline west={0.5mm}{0mm}{orange!85!white},
	title=\textcolor{orange!85!white}{\faHandPointRight[regular]\:\:Remarques},
	colback=white,
	detach title,before upper={\tcbtitle\ \quad },
	extras={enhanced, frame empty},
}

\newtcolorbox{brmqt}[1]{
	enhanced,
	breakable,
	fonttitle=\bfseries,
	frame empty,
	nobeforeafter,
	borderline west={0.5mm}{0mm}{orange!85!white},
	title=\textcolor{orange!85!white}{\faHandPointRight[regular]\:\:Remarque - #1},
	colback=white,
	extras={enhanced, frame empty},
	detach title,before upper={\tcbtitle\ \quad },
}

\newtcolorbox{consequences}[1]{
	enhanced,
	breakable,
	fonttitle=\bfseries,
	frame empty,
	nobeforeafter,
	borderline west={0.5mm}{0mm}{ProcessBlue},
	title=\textcolor{ProcessBlue}{#1},
	colback=white,
	extras={enhanced, frame empty},
	detach title,before upper={\tcbtitle\ :\quad },
}

\newtcolorbox{bmethode}{
	enhanced,
	breakable,
	fonttitle=\bfseries,
	frame empty,
	nobeforeafter,
	borderline west={0.5mm}{0mm}{blue!75!black},
	title=\textcolor{blue!75!black}{Méthode} ,
	colback=white,
	extras={enhanced, frame empty},
	attach boxed title to top left = {xshift = -3mm},
	boxed title style={size=small,colback=white,colframe=blue!75!black, arc=3mm,
	}
}

\newtcolorbox{bmethodes}{
	enhanced,
	breakable,
	fonttitle=\bfseries,
	frame empty,
	nobeforeafter,
	borderline west={0.5mm}{0mm}{blue!75!black},
	title=\textcolor{blue!75!black}{Méthodes} ,
	colback=white,
	extras={enhanced, frame empty},
	attach boxed title to top left = {xshift = -3mm},
	boxed title style={size=small,colback=blue!75!black,colframe=blue!75!black, arc=3mm,
	}
	boxed title after break style={size=small,colback=white,colframe=blue!75!black, arc=3mm,
	}
}

\newtcolorbox{bmethodet}[1]{
	enhanced,
	breakable,
	fonttitle=\bfseries,
	frame empty,
	nobeforeafter,
	borderline west={0.5mm}{0mm}{blue!75!black},
	title=\textcolor{blue!75!black}{\faTasks\:\:Méthode - #1} ,
	colback=white,
	extras={enhanced, frame empty},
	attach boxed title to top left = {xshift = -3mm},
	boxed title style={size=small,colback=white,colframe=blue!75!black, arc=3mm,
	}
}

\newtcolorbox{bdemo}{
	enhanced,
	breakable,
	fonttitle=\bfseries,
	frame empty,
	nobeforeafter,
	borderline west={0.5mm}{0mm}{violet!65!white},
	title={\faEdit[regular]\:\:Démonstration},
	%title after break=Démonstration (suite),
	colback=white,
	extras={enhanced, frame empty},
	attach boxed title to top left = {xshift = -3mm},
	boxed title style={size=small,colback=violet!65!white,colframe=violet!65!white, arc=3mm,
	}
}

\newtcolorbox{bdemos}{
	enhanced,
	breakable,
	fonttitle=\bfseries,
	frame empty,
	nobeforeafter,
	borderline west={0.5mm}{0mm}{violet!65!white},
	title={\faEdit[regular]\:\:Démonstrations},
	%title after break=Démonstration (suite),
	colback=white,
	extras={enhanced, frame empty},
	attach boxed title to top left = {xshift = -3mm},
	boxed title style={size=small,colback=violet!65!white,colframe=violet!65!white, arc=3mm,
	}
}

\newtcolorbox{bdemot}[1]{
	enhanced,
	breakable,
	fonttitle=\bfseries,
	frame empty,
	nobeforeafter,
	borderline west={0.5mm}{0mm}{violet!65!white},
	title={\faEdit[regular]\:\:Démonstration - #1},
	%title after break=Démonstration (suite),
	colback=white,
	extras={enhanced, frame empty},
	attach boxed title to top left = {xshift = -3mm},
	boxed title style={size=small,colback=violet!65!white,colframe=violet!65!white, arc=3mm,
	}
}

\newtcolorbox{boxPerso}[2]{
	enhanced,
	breakable,
	fonttitle=\bfseries,
	frame empty,
	nobeforeafter,
	borderline west={0.5mm}{0mm}{#2},
	title=#1,
	%title after break=Exemple (suite),
	colback=white,
	extras={enhanced, frame empty},
	detach title,before upper={\tcbtitle\ :\quad },
}


\newcommand*{\Coord}[3]{% 
	\ensuremath{\vv{#1}\, 
		\begin{pmatrix} 
			#2\\ 
			#3 
\end{pmatrix}}}

\newcommand*{\CoordSV}[2]{% 
	\ensuremath{\begin{pmatrix} 
			#1\\ 
			#2 
\end{pmatrix}}}

\newcommand*{\Coordp}[3]{% 
	\ensuremath{\vv{#1}\, 
		\left\lvert 
		\begin{matrix} 
			#2\\ 
			#3 
		\end{matrix}  
		\right.% Ne pas oublier le délimiteur invisible. 
}}


%personnalisation des listes non numérotées pour les différentes boites 
\newlist{itemEx}{itemize}{2}
\setlist[itemEx,1]{label=$\bullet$, font=\color{cyan!75!black}}
\setlist[itemEx,2]{label=$\bullet$, font=\color{cyan!75!black}}

\newlist{enumEx}{enumerate}{2}
\setlist[enumEx,1]{label=\arabic*), font=\bf\color{cyan!75!black}}
\setlist[enumEx,2]{label=\alph*), font=\bf\color{cyan!75!black}}

\newlist{itemMet}{itemize}{2}
\setlist[itemMet,1]{label=$\bullet$, font=\color{blue!75!black}}
\setlist[itemMet,2]{label=$\bullet$, font=\color{blue!75!black}}

\newlist{enumMet}{enumerate}{2}
\setlist[enumMet,1]{label=\arabic*), font=\bf\color{blue!75!black}}
\setlist[enumMet,2]{label=\alph*), font=\bf\color{blue!75!black}}

\newlist{itemDef}{itemize}{2}
\setlist[itemDef,1]{label=$\bullet$, font=\color{ForestGreen}}
\setlist[itemDef,2]{label=$\bullet$, font=\color{ForestGreen}}

\newlist{enumDef}{enumerate}{2}
\setlist[enumDef,1]{label=\arabic*), font=\bf\color{ForestGreen}}
\setlist[enumDef,2]{label=\alph*), font=\bf\color{ForestGreen}}

\newlist{itemProp}{itemize}{2}
\setlist[itemProp,1]{label=$\bullet$, font=\color{red!85!white}}
\setlist[itemProp,2]{label=$\bullet$, font=\color{red!85!white}}

\newlist{enumProp}{enumerate}{2}
\setlist[enumProp,1]{label=\arabic*), font=\bf\color{red!85!white}}
\setlist[enumProp,2]{label=\alph*), font=\bf\color{red!85!white}}

\newlist{itemDem}{itemize}{2}
\setlist[itemDem,1]{label=$\bullet$, font=\color{violet!85!white}}
\setlist[itemDem,2]{label=$\bullet$, font=\color{violet!85!white}}

\newlist{enumDem}{enumerate}{2}
\setlist[enumDem,1]{label=\arabic*), font=\bf\color{violet!85!white}}
\setlist[enumDem,2]{label=\alph*), font=\bf\color{violet!85!white}}

\newlist{itemRem}{itemize}{2}
\setlist[itemRem,1]{label=\ding{43}, font=\color{orange!85!white}}
\setlist[itemRem,2]{label=\ding{43}, font=\color{orange!85!white}}

\newlist{enumRem}{enumerate}{2}
\setlist[enumRem,1]{label=\arabic*), font=\bf\color{orange!85!white}}
\setlist[enumRem,2]{label=\alph*), font=\bf\color{orange!85!white}}

% nouvelle commande pour mettre entre guillemets une expression
\newcommand{\guillemets}[1]{\og \emph{#1} \fg}


\newcommand{\corr}[1]{\fcolorbox{red}{yellow!30}{#1}}

\newtcolorbox{corrBox}{enhanced,breakable,boxrule=1pt,nobeforeafter,boxsep=0pt,left=\fboxsep,right=\fboxsep,colback=yellow!30,colframe=red!75!black}

% commandes pour mettre en couleur (et en gras) dans une expression mathématique.
\newcommand{\mVert}[1]{{\color{green!55!black}\bm{#1}}}
\newcommand{\mVertWB}[1]{{\color{green!55!black}#1}}
\newcommand{\mRouge}[1]{{\color{red!85!white}\bm{#1}}}
\newcommand{\mRougeWB}[1]{{\color{red!85!white}#1}}
\newcommand{\mBleu}[1]{{\color{blue!75!black}\bm{#1}}}
\newcommand{\mBleuWB}[1]{{\color{blue!75!black}#1}}
\newcommand{\mOrange}[1]{{\color{orange!75!white}\bm{#1}}}
\newcommand{\mOrangeWB}[1]{{\color{orange!75!white}#1}}
\newcommand{\mViolet}[1]{{\color{violet!85!white}\bm{#1}}}
\newcommand{\mVioletWB}[1]{{\color{violet!85!white}#1}}

%commandes pour les flèches de distributivité.
\newcommand{\tikzmarkBB}[2][]{ \tikz[remember picture, overlay, baseline=-0.5ex]\node[#1, left](#2){}; }
\newcommand{\flechePropV}[4][]{ \tikz[remember picture, overlay]\draw (#2)edge[bend left=90,min distance=1em, -stealth]node[midway, #1]{#4}(#3); }
\newcommand{\flechePropH}[4][]{ \tikz[remember picture, overlay]\draw (#2)edge[bend left=90, min distance=1em, -stealth]node[midway, #1]{#4}(#3); }
\newcommand{\flechePropHbas}[4][]{ \tikz[remember picture, overlay]\draw (#2)edge[bend left=-90, min distance=1em, -stealth]node[midway, #1]{#4}(#3); }


\lstdefinestyle{algo}{%
	inputencoding=utf8,
	mathescape,
	basicstyle=\footnotesize\ttfamily,
	numberstyle=\scriptsize\ttfamily\color{gray},
	showstringspaces=false,
	tabsize=2,
	numbers=left,
	framexleftmargin=8mm,
	xleftmargin=8mm,
	keepspaces=true,
	frame=lines,
	framerule=1pt,
	rulecolor=\color{green!50!black},
	backgroundcolor=\color{green!5},
	keywords = {Pour, allant, Si, alors, Sinon, Fin},    
	morecomment=[l][commentstyle]{//},
	commentstyle=\color{gray},
	classoffset=0,
	keywordstyle=\color{blue!75},
	classoffset=1,
	morekeywords = {Afficher},
	keywordstyle=\color{green!50!black},
	classoffset=0,
	morestring=[b]",
	stringstyle=\color{red!75},
	literate=
	{é}{{\'e}}{1}%
	{è}{{\`e}}{1}%
	{à}{{\`a}}{1}%
	{â}{{\^a}}{1}%%%
	{ç}{{\c{c}}}{1}%
	{œ}{{\oe}}{1}%
	{ù}{{\`u}}{1}%
	{É}{{\'E}}{1}%
	{È}{{\`E}}{1}%
	{À}{{\`A}}{1}%
	{Ç}{{\c{C}}}{1}%
	{Œ}{{\OE}}{1}%
	{Ê}{{\^E}}{1}%
	{ê}{{\^e}}{1}%
	{î}{{\^i}}{1}%
	{ï}{{\"i}}{1}%%%
	{ô}{{\^o}}{1}%
	{û}{{\^u}}{1}%
	{=}{$\leftarrow$}{1}%
	{==}{$={}$}{1}}


%  --------   Macro pour les exercices et leur numérotation  

\newcounter{exo}

\setcounter{exo}{1}

\newcommand{\exercice}[1]{\vskip 0.2cm\noindent{\bf \ Exercice \theexo {  #1 }}\addtocounter{exo}{1}}

\newcounter{rit}

\setcounter{rit}{1}

\newcommand{\rituel}[1]{\vskip 0.2cm\noindent{ \bf \ \underline{Rit. \therit {  #1 }}}\addtocounter{rit}{1}}

\newcommand{\raisedrule}[2][0em]{\leaders\hbox{\rule[#1]{1pt}{#2}}\hfill}

\newcounter{numexo}

\setcounter{numexo}{1}

\newcommand{\exerciceDS}[1]{\vskip 0.6cm\noindent\rule[+2pt]{1cm}{0.5pt}\ {\large\bf \ Exercice \thenumexo {  #1 }}\raisedrule[+2pt]{0.5pt}\null\addtocounter{numexo}{1}\par\vskip 0,4cm}

\newcommand{\bareme}[1]{\tikz\node[circle,draw,inner sep=1pt,outer sep=1pt, color=red]{\textbf{#1}};}

%  --------   Pour le langage Python

% Default fixed font does not support bold face
\DeclareFixedFont{\ttb}{T1}{txtt}{bx}{n}{12} % for bold
\DeclareFixedFont{\ttm}{T1}{txtt}{m}{n}{12}  % for normal

% Custom colors
\usepackage{color}
\definecolor{deepblue}{rgb}{0,0,0.5}
\definecolor{deepred}{rgb}{0.6,0,0}
\definecolor{deepgreen}{rgb}{0,0.5,0}

\usepackage{listings}

% Python style for highlighting
\newcommand\pythonstyle{\lstset{
		language=Python,
		basicstyle=\ttm,
		otherkeywords={self},             % Add keywords here
		keywordstyle=\ttb\color{deepblue},
		emph={MyClass,__init__},          % Custom highlighting
		emphstyle=\ttb\color{deepred},    % Custom highlighting style
		stringstyle=\color{deepgreen},
		frame=tb,                         % Any extra options here
		showstringspaces=false            % 
}}


% Python environment
\lstnewenvironment{python}[1][]
{
	\pythonstyle
	\lstset{#1}
}
{}

% Python for external files
\newcommand\pythonexternal[2][]{{
		\pythonstyle
		\lstinputlisting[#1]{#2}}}

% Python for inline
\newcommand\pythoninline[1]{{\pythonstyle\lstinline!#1!}}


% \begin{document}
	
	% \section{``In-text'' listing highlighting}
	
	% \begin{python}
		% class MyClass(Yourclass):
		%     def __init__(self, my, yours):
		%         bla = '5 1 2 3 4'
		%         print bla
		% \end{python}
	
	% \section{External listing highlighting}
	
	% \pythonexternal{demo.py}
	
	% \section{Inline highlighting}
	
	% Definition \pythoninline{class MyClass} means \dots
	
	% \end{document}

\definecolor{codegreen}{rgb}{0,0.6,0}
\definecolor{codegray}{rgb}{0.5,0.5,0.5}
\definecolor{codepurple}{rgb}{0.58,0,0.82}
\definecolor{backcolour}{rgb}{0.95,0.95,0.92}

\lstdefinestyle{myStyle}{
	backgroundcolor=\color{backcolour},   
	commentstyle=\color{codegreen},
	keywordstyle=\color{magenta},
	numberstyle=\tiny\color{codegray},
	stringstyle=\color{codepurple},
	basicstyle=\ttfamily\footnotesize,
	breakatwhitespace=false,         
	breaklines=true,                 
	captionpos=b,                    
	keepspaces=true,                 
	numbers=left,                    
	numbersep=5pt,                  
	showspaces=false,                
	showstringspaces=false,
	showtabs=false,                  
	tabsize=2
}

\lstset{style=myStyle}

\newcommand{\normeB}[1]{\left\Vert \vv{#1}\right\Vert}

\newcommand\grasmaths[1]{\symbf{#1}}
\newcommand\sfbmath[1]{\symbf{\symsfup{#1}}}
\definecolor{marron}{rgb}{0.64,0.16,0.16}
\definecolor{titrebleu}{rgb}{0,0,0.6}

\usepackage{tabularx}
\usepackage{wasysym}
% ===========================tableaux de base
\newcolumntype{R}[1]{>{\raggedleft\arraybackslash }b{#1}} 
\newcolumntype{L}[1]{>{\raggedright\arraybackslash }b{#1}}
\newcolumntype{C}[1]{>{\centering\arraybackslash }b{#1}}
\newcolumntype{M}[1]{>{\centering\arraybackslash }m{#1}} %doublecentrée
\renewcommand{\tabularxcolumn}[1]{m{#1}}
\newcolumntype{Y}{>{\centering\arraybackslash }X}

%=============================boîtes de cours
\usepackage{fontawesome5}
\newcommand\fontbox{\sffamily}
%\newcommand\titrebox[1]{\textbf{#1}}
\tcbuselibrary{breakable}
\tcbset{bdefaut/.style={%
		fonttitle=\fontbox,%
		enhanced,breakable,center,boxrule=1.25pt,drop shadow=black!33!white,%
		sharp corners=uphill,arc=12pt,before skip=12pt,after skip=12pt,%
		left=5pt,leftupper=5pt,top=1mm,bottom=1mm,%
		before upper=\leavevmode
	}
}
\NewDocumentCommand\bbnewbox{ m m m m O{#4} !O{} }{%nom + titre + icone + bordure/titre/fondtitre + couleur fond + options
	\DeclareTColorBox{#1}{ O{} }{%
		bdefaut,%
		colframe=#4,coltitle=#4!25!Black,colbacktitle=#4!35,colback=#5!10,%
		title={#3\:\:\titrebox{#2##1}},%
		#6
	}%
}

\bbnewbox{bdefi}{Définition}{\faIcon[regular]{comment-dots}}{ForestGreen}
\bbnewbox{bcor}{Corollaire}{\faIcon[regular]{comment-dots}}{ForestGreen}
\bbnewbox{bdefis}{Définitions}{\faIcon[regular]{comment-dots}}{ForestGreen}
\bbnewbox{bhumour}{Humour}{\faIcon[regular]{laugh-wink}}{Yellow}
%\bbnewbox{bprop}{Propriété}{\faCog}{Salmon}
\bbnewbox{bprop}{Propriété}{\faCog}{red!85!white}
\bbnewbox{bprops}{Propriétés}{\faCog}{Salmon}
\bbnewbox{bcscq}{Conséquence}{\faArrowAltCircleRight[regular]}{Salmon}
\bbnewbox{bthm}{Théorème}{\faBullhorn}{FireBrick}
\bbnewbox{blog}{Logiciel}{\faDev}{CornflowerBlue}[white]
\bbnewbox{bhistoire}{Point histoire}{\faUniversity}{Chocolate}[white]

%\bbnewbox{brmq}{Remarque}{\faHandPointRight[regular]}{MidnightBlue}[white]
\bbnewbox{bexercice}{Exercice}{\faHubspot}{MidnightBlue}
%\bbnewbox{bexemple}{Exemple}{\faHubspot}{MediumTurquoise}
\bbnewbox{bvoc}{Vocabulaire}{\faList*[regular]}{YellowGreen}
\bbnewbox{bnota}{Notation}{\faList*}{YellowGreen}
\bbnewbox{bregle}{Règle}{\faList*}{YellowGreen}
\bbnewbox{bcadre}{Cadre}{\faInfoCircle}{Peru}
\bbnewbox{bidee}{Idée}{\faInfo}{LightBlue}
\bbnewbox{bintro}{Intro}{\faMarker}{LightSalmon}
\bbnewbox{billustr}{Illustration}{\faChartLine}{Crimson}[White]
\bbnewbox{bvideo}{}{\faIcon[regular]{file-video}}{teal}[White]
\bbnewbox{bbilan}{Bilan}{\faForward}{Purple}
%\bbnewbox{bmethode}{Méthode}{\faTasks}{Coral}
%\bbnewbox{bmethode}{Méthode}{\faTasks}{blue!75!black}
\bbnewbox{btuto}{Tuto}{\faTerminal}{Coral}[White]
\bbnewbox{bmanip}{Manipulation}{\faCodeBranch}{DarkCyan}[White]
\bbnewbox{bcalco}{Calculatrice}{\faCalculator}{MediumVioletRed}[White]
\bbnewbox{boutil}{Outil}{\faTools}{DeepSkyBlue}[White]
\bbnewbox{basavoir}{À savoir}{\faHeart}{Red}
\bbnewbox{bbloctd}{TD}{\faHourglassStart}{LightSeaGreen}[White]
\bbnewbox{bblocdm}{À vous de jouer}{\faHourglassStart}{LightSeaGreen}[White]
%\bbnewbox{brappel}{Rappel}{\faHandPointUp[regular]}{LightSalmon}
\bbnewbox{brappel}{Rappel}{\faHandPointUp[regular]}{Peru}
\bbnewbox{bresol}{Résolution}{\faPuzzlePiece}{LightSeaGreen}[White]
\bbnewbox{battention}{Attention}{\faExclamationTriangle}{Red}
\bbnewbox{bcode}{Code}{\faCode}{Crimson}[White]
\bbnewbox{bpython}{Python}{\faPython}{SeaGreen}[White]

\newcommand{\grasInDef}[1]{\textcolor{ForestGreen!55!black}{\textbf{#1}}}
\newcommand{\grasInProp}[1]{\textcolor{red!85!white}{\textbf{#1}}}
\newcommand{\grasInEx}[1]{\textcolor{MediumTurquoise!55!black}{\textbf{#1}}}
\newcommand{\grasInRem}[1]{\textcolor{orange!65!white}{\textbf{#1}}}
\newcommand{\grasInMet}[1]{\textcolor{blue!75!black}{\textbf{#1}}}
\newcommand{\grasInDem}[1]{\textcolor{violet!85!white}{\textbf{#1}}}

%\newtcolorbox{bmethodet}[1][]{%
%	bdefaut,%
%	colframe=blue!75!black,coltitle=white,colbacktitle=blue!65!black,colback=blue!75!black!15!white,
%	title={\faTasks\:\:\titrebox{{Méthode - #1}}}
%}

%outil
\newtcolorbox{boutilcalc}[1][]{%
	bdefaut,%
	colframe=DeepSkyBlue,coltitle=black,colbacktitle=DeepSkyBlue!35,colback=White,
	title=\textcolor{DeepSkyBlue!25!Black}{\faFileVideo[regular]\:\:\titrebox{Tuto#1 Vidéo}}
}

%loglogo
\newtcolorbox{bloglogo}[2][]{%
	bdefaut,%
	colframe=CornflowerBlue,coltitle=black,colbacktitle=CornflowerBlue!35,colback=white,
	title=\textcolor{CornflowerBlue!25!Black}{\faDev\:\:\titrebox{Logiciel#1}},
	overlay={\begin{tcbclipinterior}
			\node[opacity=0.33,rotate=-45] at ($(interior.south east)+(-10pt,10pt)$){\includegraphics[height=24pt]{#2}};
	\end{tcbclipinterior}}
}

%demo
%\newtcolorbox{bdemo}[1][]{%
%	bdefaut,%
%	colframe=Violet,coltitle=black,colbacktitle=Violet!35,colback=Violet!15,
%	title=\textcolor{Violet!25!Black}{\faEdit[regular]\:\:\titrebox{Démonstration#1}},
%	overlay={\begin{tcbclipinterior}%
%			\node[] at ($(interior.south east)+(-10pt,10pt)$){\large\textcolor{Violet}{\textbf{\CheckedBox}}};%
%	\end{tcbclipinterior}}
%}


%\newtcolorbox{bdemo}[1][]{%
%	bdefaut,%
%	colframe=violet!85!white,coltitle=black,colbacktitle=Violet!35,colback=Violet!15,
%	title=\textcolor{Violet!25!Black}{\faEdit[regular]\:\:\titrebox{Démonstration#1}},
%	overlay={\begin{tcbclipinterior}%
%			\node[] at ($(interior.south east)+(-10pt,10pt)$){\large\textcolor{Violet}{\textbf{\CheckedBox}}};%
%	\end{tcbclipinterior}}
%}

%demoblanc
\newtcolorbox{bdemoblanc}[1][]{%
	bdefaut,%
	colframe=Violet,coltitle=black,colbacktitle=Violet!35,colback=White,
	title=\textcolor{Violet!25!Black}{\faEdit[regular]\:\:\titrebox{Démonstration#1}},
	overlay={\begin{tcbclipinterior}%
			\node[] at ($(interior.south east)+(-10pt,10pt)$){\large\textcolor{Violet}{\textbf{\CheckedBox}}};%
	\end{tcbclipinterior}}
}

%algo
\newtcolorbox{balgo}[1][]{%
	bdefaut,%
	colframe=OrangeRed,coltitle=black,colbacktitle=OrangeRed!35,colback=White,
	title=\textcolor{OrangeRed!25!Black}{\faFileCode\:\:\titrebox{Algorithme#1}},
	overlay={\begin{tcbclipinterior}%
			\node[opacity=0.15,rotate=-15] at ($(interior.south east)+(-8pt,12pt)$){\includegraphics[height=24pt]{algorithmvec}};%
	\end{tcbclipinterior}}
}


%aide
\newtcolorbox{baide}{%
	bdefaut,%
	colframe=titrebleu!35,coltitle=titrebleu!50!Black,colbacktitle=titrebleu!15,colback=White,
	title=\textcolor{titrebleu!25!Black}{\faBacon\:\:\titrebox{Niveaux de difficulté}},
	overlay={\begin{tcbclipinterior}%
			\node[opacity=0.33,rotate=45,scale=0.33] at ($(interior.south east)+(-6pt,6pt)$){\textcolor{titrebleu!50!Black}{\footnotesize \bfseries \sffamily \textcopyright{} Dairon11}};%
	\end{tcbclipinterior}}
}

\newtcolorbox{bcalcologo}[2][]{%
	bdefaut,%
	colframe=MediumVioletRed,coltitle=black,colbacktitle=MediumVioletRed!35,colback=White,
	title=\textcolor{MediumVioletRed!25!Black}{\faCalculator\:\:\titrebox{Calculatrice#1}},
	overlay={\begin{tcbclipinterior}%
			\node[opacity=0.33,rotate=-45] at ($(interior.south east)+(-10pt,10pt)$){\includegraphics[height=48pt]{#2}};%
	\end{tcbclipinterior}}
}

\definecolor{orangescratch}{HTML}{F8A31F}
\newtcolorbox{bscratch}[1][]{%
	bdefaut,%
	colframe=orangescratch,coltitle=black,colbacktitle=orangescratch!35,colback=White,
	title=\textcolor{orange!25!Black}{\includegraphics[height=1em]{scratch}\:\:\titrebox{Scratch#1}},
	overlay={\begin{tcbclipinterior}%
			\node[opacity=0.33,rotate=45] at ($(interior.south east)+(-12pt,12pt)$){\includegraphics[height=12pt]{scratchban}};%
	\end{tcbclipinterior}}
}

%reponse
\newtcolorbox{breponse}[1][]{%
	bdefaut,%
	colframe=Thistle,coltitle=black,colbacktitle=Thistle!35,colback=White,
	before upper = \vspace{-0.05cm},
	title=\textcolor{Thistle!25!Black}{\faIcon{pen-alt}\:\:\titrebox{Réponse#1}}
}

\usepackage{dashundergaps}
\newcommand\underdash[1][black]{% <- règle la couleur du texte
	\bgroup
	\ifdim\ULdepth=\maxdimen\settodepth\ULdepth{(j}\advance\ULdepth.4pt\fi
	\markoverwith{\kern0.2em% <- règle l'espacement avant un pointillé
		\vtop{\kern0.5ex% <- règle l'altitude des pointillés
			{\color{black}% <- règle la couleur des pointillés
				\hrule width.4em% <- règle la largeur des pointillés
		}}\kern0.2em% <- règle l'espacement après un pointillé
	}\color{#1}\ULon}

\newcommand\pointille[1][1cm]{\makebox[#1]{\dotfill}}



\newcommand{\va}[1]{\lvert \ #1 \ \rvert}
\newcommand{\vaf}[1]{\left\lvert \ #1\  \right\rvert}

\begin{document}
\graphicspath{{images/}}
\setcounter{chapter}{5}
\chapter[Correction des exercices]{Suites numériques \\ correction des exercices}
\thispagestyle{fancy}

\section*{généralité sur les suites numériques}

\subsection*{Ex 40 p 32}

\begin{enumMet}
\item $u_{0}:$ impossible à calculer $;\quad u_{1}=\sqrt{1-1}+2 \times 1=2$\quad $u_{2}=\sqrt{2-1}+2 \times 2=5 ;$\\
$u_{3}=\sqrt{3-1}+2 \times 3=\sqrt{2}+6$\quad
$u_{4}=\sqrt{4-1}+2 \times 4=\sqrt{3}+8 ;\quad$\\ $u_{10}=\sqrt{10-1}+2 \times 10=23$
    \item $u_{0}=\dfrac{5 \times 0-3}{2 \times 0-2}=\dfrac{-3}{-2}=\dfrac{3}{2} ;\quad u_{1}:$ impossible à calculer\quad
$u_{2}=\dfrac{5 \times 2-3}{2 \times 2-2}=\dfrac{7}{2},$\\ $u_{3}=\dfrac{5 \times 3-3}{2 \times 3-2}=\dfrac{12}{4}=3$\quad
$u_{4}=\dfrac{5 \times 4-3}{2 \times 4-2}=\dfrac{17}{6},\quad u_{10}=\dfrac{5 \times 10-3}{2 \times 10-2}=\dfrac{47}{18}$
    \item $u_{0}=\cos \left(0 \times \dfrac{\pi}{2}\right)=1,\quad u_{1}=\cos \left(1 \times \dfrac{\pi}{2}\right)=0$\quad
$u_{2}=\cos \left(2 \times \dfrac{\pi}{2}\right)=-1,$\\
$u_{3}=\cos \left(3 \times \dfrac{\pi}{2}\right)=0$ \quad
$u_{4}=\cos \left(4 \times \dfrac{\pi}{2}\right)=1, \quad u_{10}=\cos \left(10 \times \dfrac{\pi}{2}\right)=-1$
    \item $u_{0}=\left(\dfrac{1}{2}\right)^{0}=1 \quad u_{1}=\left(\dfrac{1}{2}\right)^{1}=\dfrac{1}{2}$\quad
$u_{2}=\left(\dfrac{1}{2}\right)^{2}=\dfrac{1}{4}$\\ $u_{3}=\left(\dfrac{1}{2}\right)^{3}=\dfrac{1}{8}$\quad
$u_{4}=\left(\dfrac{1}{2}\right)^{4}=\dfrac{1}{16}\quad u_{10}=\left(\dfrac{1}{2}\right)^{10}=\dfrac{1}{1024}$
\end{enumMet}

\subsection*{Ex 43 p 32}

\begin{enumMet}
    \item $u_{n}=6 n+8$\\
    $u_{n-1}=6(n-1)+8=6 n+2$\\
    $u_{n+1}=6(n+1)+8=6 n+14$
    \item $u_{n}=n^{2}-2 n+8$\\
    $u_{n-1}=(n-1)^{2}-2(n-1)+8=n^{2}-2 n+1-2 n+2+8=n^{2}-4 n+11$\\
    $u_{n+1}=(n+1)^{2}-2(n+1)+8=n^{2}+2 n+1-2 n-2+8=n^{2}+7$
    \item $u_{n}=\dfrac{n(n+1)}{n+2}$\\
    $u_{n-1}=\dfrac{(n-1)(n-1+1)}{n-1+2}=\dfrac{(n-1) n}{n+1}$\\
    $u_{n+1}=\dfrac{(n+1)(n+1+1)}{n+1+2}=\dfrac{(n+1)(n+2)}{n+3}$\\
    \item $u_{n}=5^{n}$\\
    $u_{n-1}=5^{n-1}=\dfrac{5^{n}}{5}$\\
    $u_{n+1}=5^{n+1}=5 \times 5^{n}$\\
    \item $u_{n}=\dfrac{3^{n+1}}{2^{n}}$\\
    $u_{n-1}=\dfrac{3^{n-1+1}}{2^{n-1}}=\dfrac{3^{n}}{2^{n-1}}$\\
    $u_{n-1}=\dfrac{3^{n+1+1}}{2^{n+1}}=\dfrac{3^{n+2}}{2^{n+1}}$
    \item $u_{n}=\dfrac{9 n-5}{4 n+6}$\\
    $u_{n-1}=\dfrac{9(n-1)-5}{4(n-1)+6}=\dfrac{9 n-14}{4 n+2}$\\
    $u_{n+1}=\dfrac{9(n+1)-5}{4(n+1)+6}=\dfrac{9 n+4}{4 n+10}$
    \item $u_{n}=\left(\dfrac{n^{2}}{n+1}\right)^{n+1}$\\
    $u_{n-1}=\left(\dfrac{(n-1)^{2}}{n-1+1}\right)^{n-1+1}=\left(\dfrac{(n-1)^{2}}{n}\right)^{n}=\left(\dfrac{n^{2}-2 n+1}{n}\right)^{n}$\\
    $u_{n+1}=\left(\dfrac{(n+1)^{2}}{n+1+1}\right)^{n+1+1}=\left(\dfrac{(n+1)^{2}}{n+2}\right)^{n+2}=\left(\dfrac{n^{2}+2 n+1}{n+2}\right)^{n+2}$
\end{enumMet}

\subsection*{Ex 44 p 33}

\begin{enumMet}
    \item $u_{1}=3 u_{0}-4 \times 0=3 \times 2-4 \times 0=6$\\
    $u_{2}=3 u_{1}-4 \times 1=3 \times 6-4 \times 1=14$\\
    $u_{3}=3 u_{2}-4 \times 2=3 \times 14-4 \times 2=34$
    \item $u_{1}=u_{0}^{2}+\dfrac{1}{2 \times 0+1}=0^{2}+\dfrac{1}{2 \times 0+1}=1$\\
    $u_{2}=u_{1}^{2}+\dfrac{1}{2 \times 1+1}=1^{2}+\dfrac{1}{2 \times 1+1}=1+\dfrac{1}{3}=\dfrac{4}{3}$\\
    $u_{3}=u_{2}^{2}+\dfrac{1}{2 \times 2+1}=\left(\dfrac{4}{3}\right)^{2}+\dfrac{1}{2 \times 2+1}=\dfrac{89}{45}$
    \item $u_{1}=5 u_{0}-2=5 \times 3-2=13$\\
    $u_{2}=5 u_{1}-2=5 \times 13-2=63$\\
    $u_{3}=5 u_{2}-2=5 \times 63-2=313$
    \item $u_{2}=2 u_{1}+u_{0}=2 \times 2+1=5$\\
    $u_{3}=2 u_{2}+u_{1}=2 \times 5+2=12$\\
    $u_{4}=2 u_{3}+u_{2}=2 \times 12+5=29$
    \item $u_{2}=u_{1}+0=2+0=2$\\
    $u_{3}=u_{2}+1=2+1=3$\\
    $u_{4}=u_{3}+2=3+2=5$
\end{enumMet}

\subsection*{Ex 22 p 31}

\begin{enumMet}
    \item $u_{0}=3 \times 0+7=7 ;\quad u_{1}=3 \times 1+7=10 ;\quad u_{2}=3 \times 2+7=13$\\
    $u_{3}=3 \times 3+7=16 ;\quad  u_{10}=3 \times 10+7=37$
    \item $u_{0}=\dfrac{0}{0+1}+2=2$\quad $u_{1}=\dfrac{1}{1+1}+2=\dfrac{5}{2}$; \quad $u_{2}=\dfrac{2}{2+1}+2=\dfrac{8}{3}$\\
    $u_{3}=\dfrac{3}{3+1}+2=\dfrac{11}{4} \quad u_{10}=\dfrac{10}{10+1}+2=\dfrac{32}{11}$
    \item $u_{0}=0^{2}-7 \times 0+2=2 ;\quad u_{1}=1^{2}-7 \times 1+2=-4 \quad u_{2}=2^{2}-7 \times 2+2=-8 ;$ \\
    $u_{3}=3^{2}-7 \times 3+2=-10 \quad u_{10}=10^{2}-7 \times 10+2=32 $
    \item $u_{0}=2^{0}-3=-2 ;\quad u_{1}=2^{1}-3=-1 ;\quad u_{2}=2^{2}-3=1$\\
    $u_{3}=2^{3}-3=5 ;\quad  u_{10}=2^{10}-3=1021$
    \item $u_{0}=(-1)^{0}-0=1 ;\quad u_{1}=(-1)^{1}-1=-2 ;\quad u_{2}=(-1)^{2}-2=-1$\\
    $u_{3}=(-1)^{3}-3=-4 ;\quad u_{10}=(-1)^{10}-10=-9$
    \item $u_{0}=\dfrac{3}{4}-\dfrac{1}{4} \times 0=\dfrac{3}{4}$;\quad  $u_{1}=\dfrac{3}{4}-\dfrac{1}{4} \times 1=\dfrac{1}{2}$;\quad $u_{2}=\dfrac{3}{4}-\dfrac{1}{4} \times 2=\dfrac{1}{4}$\\
    $u_{3}=\dfrac{3}{4}-\dfrac{1}{4} \times 3=0;\quad u_{10}=\dfrac{3}{4}-\dfrac{1}{4} \times 10=-\dfrac{7}{4}$
\end{enumMet}

\subsection*{Ex 32 p 32}

\begin{enumMet}
    \item $u_{0}=7 \times 0+1=1 ;\quad u_{1}=7 \times 1+1=8 ;\quad u_{2}=7 \times 2+1=15$\\
    $u_{3}=7 \times 3+1=22 ;\quad u_{4}=7 \times 4+1=29$
    \item $v_{0}=0^{2}-3 \times 0+1=1 ;\quad v_{1}=1^{2}-3 \times 1+1=-1$\quad
    $v_{2}=2^{2}-3 \times 2+1=-1 ;$\\
    $v_{3}=3^{2}-3 \times 3+1=1$\quad 
    $v_{4}=4^{2}-3 \times 4+1=5$
    \item $w_{0}=\dfrac{1}{2 \times 0-1}=\dfrac{1}{-1}=-1 ;\quad w_{1}=\dfrac{1}{2 \times 1-1}=\dfrac{1}{1}=1$\quad
    $w_{2}=\dfrac{1}{2 \times 2-1}=\dfrac{1}{3},$\\
    $w_{3}=\dfrac{1}{2 \times 3-1}=\dfrac{1}{5}$\quad
    $w_{4}=\dfrac{1}{2 \times 4-1}=\dfrac{1}{7}$
    \item $t_{0}=\sqrt{5 \times 0+1}=\sqrt{1}=1 ;\quad t_{1}=\sqrt{5 \times 1+1}=\sqrt{6}$\quad
    $t_{2}=\sqrt{5 \times 2+1}=\sqrt{11} ;$\\
    $ t_{3}=\sqrt{5 \times 3+1}=\sqrt{16}=4$\quad
    $t_{4}=\sqrt{5 \times 4+1}=\sqrt{21}$
\end{enumMet}

\subsection*{Ex 33 p 32}

\begin{enumMet}
    \item $u_{0}=2 ;\quad\quad u_{1}=2 u_{0}+1=2 \times 2+1=5$\\
    $u_{2}=2 u_{1}+1=2 \times 5+1=11 ; \quad\quad u_{3}=2 u_{2}+1=2 \times 11+1=23$\\
    $u_{4}=2 u_{3}+1=2 \times 23+1=47$
    \item $u_{0}=1 ; \quad\quad u_{1}=u_{0}^{2}+1=1^{2}+1=2$\\
    $u_{2}=u_{1}^{2}+1=2^{2}+1=5 ;\quad\quad u_{3}=u_{2}^{2}+1=5^{2}+1=26$\\
    $u_{4}=u_{3}^{2}+1=26^{2}+1=677$
    \item $u_{1}=1 ;\quad \quad u_{2}=\dfrac{1}{u_{1}+2}=\dfrac{1}{1+2}=\dfrac{1}{3}$ \\ $u_{3}=\dfrac{1}{u_{2}+2}=\dfrac{1}{\dfrac{1}{3}+2}=\dfrac{3}{7} \quad \quad u_{4}=\dfrac{1}{u_{3}+2}=\dfrac{1}{\dfrac{3}{7}+2}=\dfrac{7}{17} \\ u_{5}=\dfrac{1}{u_{4}+2}=\dfrac{1}{\dfrac{7}{17}+2}=\dfrac{17}{41}$
    \item $u_{0}=1$ ;\quad\quad $u_{1}=\sqrt{u_{0}^{2}+1}=\sqrt{1^{2}+1}=\sqrt{2}$\\
    $u_{2}=\sqrt{u_{1}^{2}+1}=\sqrt{(\sqrt{2})^{2}+1}=\sqrt{3} \quad \quad u_{3}=\sqrt{u_{2}^{2}+1}=\sqrt{(\sqrt{3})^{2}+1}=\sqrt{4}=2$\\
    $u_{4}=\sqrt{u_{3}^{2}+1}=\sqrt{2^{2}+1}=\sqrt{5}$
\end{enumMet}

\subsection*{Ex 38 p 32}
\grasInMet{Soit $\left(u_{n}\right)$ la suite définie pour tout entier naturel $n$ par $u_{n}=\dfrac{1}{4 n+3}$. Déterminer son sens de variation.}\\

Pour tout entier $n \in \mathbb{N}, u_{n+1}=\dfrac{1}{4(n+1)+3}=\dfrac{1}{4 n+7}$\\
Pour tout entier $n \in \mathbb{N}, u_{n+1}-u_{n}=\dfrac{1}{4 n+7}-\dfrac{1}{4 n+3}$ soit\\ $u_{n+1}-u_{n}=\dfrac{4 n+3}{(4 n+7)(4 n+3)}-\dfrac{4 n+7}{(4 n+3)(4 n+7)}$\\ $u_{n+1}-u_{n}=\dfrac{-4}{(4 n+7)(4 n+3)}$
Comme $n \geqslant 0,4 n+3 \geqslant 0$ et $4 n+7 \geqslant 0$,
Donc, pour tout entier $n \in \mathbb{N}, u_{n+1}-u_{n} \leqslant 0$. La suite est décroissante.

\subsection*{Ex 49 p 33}
\grasInMet{Dans chaque cas, déterminer le sens de variation de la suite $\left(u_{n}\right)$ définie par:}
\begin{enumMet}
    \item \grasInMet{$u_{n}=\dfrac{n+1}{n+2}$ pour tout $n \geqslant 0$}\\
    Pour tout entier naturel $n$ :
    $$
    \begin{aligned}
    &u_{n+1}=\dfrac{n+1+1}{n+1+2}=\dfrac{n+2}{n+3} \\
    &u_{n+1}-u_{n}=\dfrac{n+2}{n+3}-\dfrac{n+1}{n+2}=\dfrac{(n+2)^{2}-(n+1)(n+3)}{(n+2)(n+3)} \\
    &u_{n+1}-u_{n}=\dfrac{n^{2}+4 n+4-\left(n^{2}+4 n+3\right)}{(n+2)(n+3)} \\
    &u_{n+1}-u_{n}=\dfrac{1}{(n+2)(n+3)}
    \end{aligned}
    $$
    Comme $n \geq 0$ alors $n+2 \geq 2$ et $n+3 \geq 3, u_{n+1}-u_{n}>0$
    La suite est donc croissante.
    \item \grasInMet{$u_{n}=\dfrac{3^{n}}{n}$ pour tout $n \geqslant 1$}\\
    On a $u_{n}>0$ pour tout entier naturel, $n \geqslant 1$,\\
    $u_{n+1}=\dfrac{3^{n+1}}{n+1}$ et
    $u_{n+1}-u_{n}=\dfrac{3^{n+1}}{n+1}-\dfrac{3^{n}}{n}=\dfrac{3^{n+1} \times n}{n(n+1)}-\dfrac{3^{n}(n+1)}{n(n+1)}$\\
    $u_{n+1}-u_{n}=\dfrac{3^{n} \times 3 \times n-3^{n}(n+1)}{n(n+1)}=\dfrac{3^{n}(3 n-(n+1))}{n(n+1)}$\\
    $u_{n+1}-u_{n}=\dfrac{\left.3^{n}(3 n-n-1)\right)}{n(n+1)}=$
    $\dfrac{\left.3^{n}(2 n-1)\right)}{n(n+1)}$
    Or $3^{n}>0$ et comme $n \geq 1, n(n+1)>0$. Donc $u_{n+1}-u_{n}$ est du signe de $2 n-1$.
    Comme $n>1$, on a donc $2 n>2$ et donc $2 n-1>1>0$. D'où $u_{n+1}-u_{n}>0$. La suite est donc croissante.
    \item \grasInMet{$u_{n}=n^{2}-3 n+12$ pour tout $n \geqslant 0$}\\
    Pour tout entier naturel $n$ :\\
    $u_{n+1}=(n+1)^{2}-3(n+1)+12=n^{2}+2 n+1-3 n-3+12=n^{2}-n+10$.
    $u_{n+1}-u_{n}=n^{2}-n+10-\left(n^{2}-3 n+12\right)=2 n-2$
    On a $2 n-2>0$ si $2 n>2 \Longleftrightarrow n>1$. La suite est donc croissante à partir de $n=1$.
    \item \grasInMet{$u_{n}=\dfrac{1}{n}-\dfrac{1}{n+1}$ pour tout $n \geqslant 1$}\\
    Pour tout entier naturel $n \geqslant 1$ :
    $$
    u_{n}=\dfrac{n+1}{n(n+1)}-\dfrac{n}{n(n+1)}=\dfrac{1}{n(n+1)} \text {. }
    $$
    Les termes de la suite sont donc strictement positifs.
    $u_{n+1}=\dfrac{1}{(n+1)(n+2)}$.
    D'où $\dfrac{u_{n+1}}{u_{n}}=\dfrac{1}{(n+1)(\underset{n}{n}+2)} \times n(n+1)=\dfrac{n}{n+2}$.
    Or $n+2>n$ donc $\dfrac{n}{n+2}<1$.
    Par conséquent, pour tout entier $n \geqslant 1$, on a $u_{n}>0$ et $n \geqslant 1$, on a $u_{n}>0$ et $\dfrac{u_{n+1}}{u_{n}}<1$.
    La suite est donc décroissante.
\end{enumMet}

\subsection*{Ex 39 p 32}
\grasInMet{Étudier la monotonie des suites suivantes}
\begin{enumMet}
    \item \grasInMet{Pour tout entier naturel $n, u_{n}=2,5^{n}$.}\\
    $u_{n+1}-u_n=2,5^{n+1}-2,5^n=2,5^n(2,5-1)=1,5\times2,5^n>0$. La suite est croissante.
    \item \grasInMet{Pour tout entier naturel $n, v_{n}=0,4^{n}$.}\\
    $v_{n+1}-v_n=0,4^{n+1}-0,4^n=0,4^n(0,4-1)=-0,6\times0,4^n<0$. La suite est décroissante.
    \item \grasInMet{Pour tout entier naturel $n, w_{n}=(-1,8)^{n}$.}
    $w_{n+1}-w_n=(-1,8)^{n+1}-(-1,8)^n=(-1,8)^n(-1,8-1)=-2,8\times(-1,8)^n$.\\
    Si n est pair, $w_{n+1}-w_n<0$ et Si n est impair, $w_{n+1}-w_n>0$. La suite n'est donc pas monotone.
    \item \grasInMet{Pour tout entier naturel $n, t_{n}=3 n+5$.}\\
    Pour tout entier $t_{n+1}=3(n+1)+5=3 n+8$
    Pour tout entier $t_{n+1}-t_{n}=3 n+8-(3 n+5)=3>0$
    La suite $\left(t_{n}\right)$ est croissante.\\
\end{enumMet}

\section*{Suites arithmétiques}


\subsection*{Ex 55 p 34}

\begin{enumMet}
    \item $u_{n}=\dfrac{n+5}{n+1}$\\
    On calcule les premiers termes:\\
    $u_{0}=\dfrac{0+5}{0+1}=5\quad  u_{1}=\dfrac{1+5}{1+1}=3 \text {\quad et\quad  } u_{2}=\dfrac{2+5}{2+1}=\dfrac{7}{3}$\\
    On a $u_{2}-u_{1}=-\dfrac{2}{3}$ et $u_{1}-u_{0}=-2$\\
    La suite n'est donc pas arithmétique.
    \item $u_{n}=\dfrac{-3 n+5}{8}$\\
    On calcule les premiers termes:\\
    $u_{0}=\dfrac{5}{8}\quad u_{1}=\dfrac{-3+5}{8}=\dfrac{2}{8}=\dfrac{1}{4}$ \quad et\quad 
    $u_{2}=\dfrac{-3 \times 2+5}{8}=-\dfrac{1}{8}$\\
    On a $u_{2}-u_{1}=-\dfrac{3}{8}$ et $u_{1}-u_{0}=-\dfrac{3}{8}$, la suite semble arithmétique.\\
    Il faut vérifier pour tout entier naturel $n$.\\
    $u_{n+1}=\dfrac{-3(n+1)+5}{8}=\dfrac{-3 n+2}{8}$\\
    D'où $u_{n+1}-u_{n}=\dfrac{-3 n+2}{8}-\dfrac{-3 n+5}{8}=-\dfrac{3}{8}$\\
    La suite est arithmétique de raison $r=-\dfrac{3}{8}$ et de premier terme $u_{0}=\dfrac{5}{8}$.
    \item $u_{n}=\dfrac{n^{2}+4 n+3}{n+3}$\\
    On calcule les premiers termes:\\
    $u_{0}=\dfrac{3}{3}=1\quad u_{1}=\dfrac{1^{2}+4 \times 1+3}{1+3}=\dfrac{8}{4}=2$\quad et \quad   $u_2=\dfrac{2^{2}+4 \times 2+3}{2+3}=\dfrac{15}{5}=3$\\
    On a $u_{2}-u_{1}=1$ et $u_{1}-u_{0}=1,$ la suite semble arithmétique. II faut vérifier pour tout entier naturel $n$.\\
    $u_{n+1}=\dfrac{(n+1)^{2}+4(n+1)+3}{n+1+3}$\\
    $u_{n+1}=\dfrac{n^{2}+2 n+1+4 n+4+3}{n+4}=\dfrac{n^{2}+6 n+8}{n+4}$\\
    D'où $u_{n+1}-u_{n}=\dfrac{n^{2}+6 n+8}{n+4}-\dfrac{n^{2}+4 n+3}{n+3}$\\
    $=\dfrac{\left(n^{2}+6 n+8\right)(n+3)-\left(n^{2}+4 n+3\right)(n+4)}{(n+4)(n+3)}$\\
    $=\dfrac{n^{3}+3 n^{2}+6 n^{2}+18 n+8 n+24-\left(n^{3}+4 n^{2}+4 n^{2}+16 n+3 n+12\right)}{\left(n^{2}+3 n+4 n+12\right)}$\\
    $=\dfrac{n^{2}+7 n+12}{n^{2}+7 n+12}=1$\\
    La suite est arithmétique de raison $r=1$ et de premier terme $u_{0}=1$. 
    \item $u_{n}=\dfrac{n^{2}+3}{n+5}$\\
    On calcule les premiers termes $u_{0}=\dfrac{3}{5},\quad u_{1}=\dfrac{4}{6}=\dfrac{2}{3}$ \quad et\quad  $u_{2}=\dfrac{7}{7}=1$\\
    On a $u_{2}-u_{1}=\dfrac{1}{3}$ et $u_{1}-u_{0}=\dfrac{1}{15}$. La suite n'est donc pas arithmétique.
\end{enumMet}

\subsection*{Ex 26 p 31}
\grasInMet{Les suites suivantes sont arithmétiques. Exprimer $u_{n+1}$ en fonction de $u_{n}$ puis calculer les cinq premiers termes.}
\begin{enumMet}
    \item \grasInMet{$u_{0}=2$ et $r=3$}\\
    Pour tout $n \in \mathbb{N}, u_{n+1}=u_{n}+3$
    $u_{0}=2$
    $$
    \begin{aligned}
    &u_{1}=u_{0}+3=2+3=5 \\
    &u_{2}=u_{1}+3=5+3=8 \\
    &u_{3}=u_{2}+3=8+3=11 \\
    &u_{4}=u_{3}+3=11+3=14
    \end{aligned}
    $$
    \item \grasInMet{$u_{0}=4$ et $r=-3$}\\
    Pour tout $n \in \mathbb{N}, u_{n+1}=u_{n}-3$
    $$
    \begin{aligned}
    &u_{0}=4 \\
    &u_{1}=u_{0}-3=4-3=1 \\
    &u_{2}=u_{1}-3=1-3=-2 \\
    &u_{3}=u_{2}-3=-2-3=-5 \\
    &u_{4}=u_{3}-3=-5-3=-8
    \end{aligned}
    $$
    \item \grasInMet{$u_{1}=3$ et $r=\frac{1}{2}$}\\
    $$
    \begin{aligned}
    &\text { Pour tout entier } n \geqslant 1, u_{n+1}=u_{n}+\frac{1}{2} \\
    &u_{1}=3 \\
    &u_{2}=u_{1}+\frac{1}{2}=3+\frac{1}{2}=\frac{7}{2} \\
    &u_{3}=u_{2}+\frac{1}{2}=\frac{7}{2}+\frac{1}{2}=\frac{8}{2}=4 \\
    &u_{4}=u_{3}+\frac{1}{2}=4+\frac{9}{2}=\frac{9}{2} \\
    &u_{5}=u_{4}+\frac{1}{2}=\frac{9}{2}+\frac{1}{2}=\frac{10}{2}=5
    \end{aligned}
    $$
    \item \grasInMet{$u_{1}=\frac{3}{4}$ et $r=\frac{1}{2}$}\\
    Pour tout entier $n \geqslant 1, u_{n+1}=u_{n}+\frac{1}{2}$
    $$
    \begin{aligned}
    &u_{1}=\frac{3}{4} \\
    &u_{2}=u_{1}+\frac{1}{2}=\frac{3}{4}+\frac{1}{2}=\frac{5}{4} \\
    &u_{3}=u_{2}+\frac{1}{2}=\frac{5}{4}+\frac{1}{2}=\frac{7}{4} \\
    &u_{4}=u_{3}+\frac{1}{2}=\frac{7}{4}+\frac{1}{2}=\frac{9}{4} \\
    &u_{5}=u_{4}+\frac{1}{2}=\frac{9}{4}+\frac{1}{2}=\frac{11}{4}
    \end{aligned}
    $$
\end{enumMet}

\subsection*{Ex 27 p 31}
% \begin{enumMet}
%     \item $u_{0}=2 ;\quad\quad u_{1}=2 u_{0}+1=2 \times 2+1=5$\\
%     $u_{2}=2 u_{1}+1=2 \times 5+1=11 ; \quad\quad u_{3}=2 u_{2}+1=2 \times 11+1=23$\\
%     $u_{4}=2 u_{3}+1=2 \times 23+1=47$
%     \item $u_{0}=1 ; \quad\quad u_{1}=u_{0}^{2}+1=1^{2}+1=2$\\
%     $u_{2}=u_{1}^{2}+1=2^{2}+1=5 ;\quad\quad u_{3}=u_{2}^{2}+1=5^{2}+1=26$\\
%     $u_{4}=u_{3}^{2}+1=26^{2}+1=677$
%     \item $u_{1}=1 ;\quad \quad u_{2}=\dfrac{1}{u_{1}+2}=\dfrac{1}{1+2}=\dfrac{1}{3}$ \\ $u_{3}=\dfrac{1}{u_{2}+2}=\dfrac{1}{\dfrac{1}{3}+2}=\dfrac{3}{7} \quad \quad u_{4}=\dfrac{1}{u_{3}+2}=\dfrac{1}{\dfrac{3}{7}+2}=\dfrac{7}{17} \\ u_{5}=\dfrac{1}{u_{4}+2}=\dfrac{1}{\dfrac{7}{17}+2}=\dfrac{17}{41}$
%     \item $u_{0}=1$ ;\quad\quad $u_{1}=\sqrt{u_{0}^{2}+1}=\sqrt{1^{2}+1}=\sqrt{2}$\\
%     $u_{2}=\sqrt{u_{1}^{2}+1}=\sqrt{(\sqrt{2})^{2}+1}=\sqrt{3} \quad \quad u_{3}=\sqrt{u_{2}^{2}+1}=\sqrt{(\sqrt{3})^{2}+1}=\sqrt{4}=2$\\
%     $u_{4}=\sqrt{u_{3}^{2}+1}=\sqrt{2^{2}+1}=\sqrt{5}$
% \end{enumMet}

\begin{enumMet}
    \item Pour tout entier $n \in \mathbb{N}, u_{n}=7 n-3$\\
    Pour tout entier $n \in \mathbb{N}, u_{n+1}=7(n+1)-3=7 n+4$\\
    D'où, pour tout entier $n \in \mathbb{N}$, $u_{n+1}-u_{n}=7 n+4-(7 n-3)=7 n+4-7 n+3=7$\\
    La suite est donc arithmétique de raison $r=7$ et de premier terme $u_{0}=7\times0-3=-3$.\\
    \textbf{Remarque:} On peut également reconnaître que $u_{n}$ est de la forme $u_{0}+n r$ et déduire directement que $\left(u_{n}\right)$ est arithmétique de raison $r=7$ et de premier terme $u_{0}=-3$.
    \item Pour tout $n \in \mathbb{N}, u_{n}=n^{2}+4$\\
    On calcule les premiers termes:\\
    $u_{0}=4, u_{1}=1+4=5$ et $u_{2}=2^{2}+4=8$\\
    On a $u_{2}-u_{1}=3$ et $u_{1}-u_{0}=1$ La suite n'est donc pas arithmétique.
    \item On calcule les premiers termes:\\
    $u_{1}=2 \times 1+7=9$ et $u_{2}=2 \times 9+7=25$\\
    On a $u_{2}-u_{1}=16$ et $u_{1}-u_{0}=8$ La suite n'est pas arithmétique.
    \item On calcule les premiers termes:\\
    $u_{2}=u_{1}+1+9=-3+1-9=7$\\
    $u_{3}=u_{2}+2+9=7+2+9=18$\\
    On a $u_{3}-u_{2}=11$ et $u_{2}-u_{1}=10$\\
    La suite n'est pas arithmétique.
\end{enumMet}

\subsection*{Ex 28 p 31}

La suite étant arithmétique, on sait que le terme général peut s'écrire sous la forme $u_{n}=u_{p}+(n-p) r$. On obtient donc :
\begin{enumMet}
    \item Pour tout entier $n \in \mathbb{N}, u_{n}=u_{2}+(n-2) \times r$ donc $u_{n}=2+4(n-2)=4 n-6$
    \item Pour tout entier $n \in \mathbb{N}, u_{n}=u_{5}+(n-5) \times r$ donc $u_{n}=7-\dfrac{1}{2}(n-5)=-\dfrac{1}{2} n+\dfrac{19}{2}$
    \item Pour tout entier $n \in \mathbb{N}, n \in \mathbb{N}, u_{n}=u_{3}+(n-3) \times r$ donc
    $u_{n}=4+12(n-3)=12 n-32$
    \item Pour tout entier, $n \in \mathbb{N}, u_{n}=u_{8}+(n-8) \times r$ donc $u_{n}=\dfrac{37}{2}-\dfrac{1}{4}(n-8)=-\dfrac{1}{4} n+\dfrac{41}{2}$
\end{enumMet}

\subsection*{Ex 54 p 34}

\begin{enumMet}
    \item Comme la suite est arithmétique de premier terme $u_{7}=12$ et de raison $r=3$, on peut écrire son terme général sous la forme $u_{n}=u_{p}+(n-p) r$ soit ici $: u_{n}=u_{7}+(n-7) \times 3$ On a donc $u_{20}=12+(20-7) \times 3=12+13 \times 3=51$
    \item Comme la suite est arithmétique avec $u_{25}=17$ et de raison $r=5$, on peut écrire son terme général sous la forme $u_{n}=u_{p}+(n-p) r$ soit ici : $u_{n}=u_{25}+(n-25) \times 5$. On a donc $u_{20}=17+(20-25) \times 5=17+(-5) \times 5=-8$
    \item La suite $\left(u_{n}\right)$ est arithmétique car $u_{n+1}-u_{n}=7$. La raison est donc 7 et le premier terme $u_{0}=3$
    On peut donc écrire le terme général sous la forme $u_{n}=u_{p}+(n-p) r$ soit ici :
    $u_{n}=u_{0}+n \times 7$
    On a donc $u_{20}=3+20 \times 7=3+20 \times 7=143$
    \item La suite $\left(u_{n}\right)$ est arithmétique car $u_{n+1}-u_{n}=-4$. La raison est donc -4 et le premier terme $u_{1}=-2$
    On peut donc écrire le terme général sous la forme $u_{n}=u_{p}+(n-p) r$ soit ici :
    $u_{n}=u_{1}+(n-1) \times(-4)$
\end{enumMet}



\subsection*{Ex 56 p 34}

La suite étant arithmétique, on sait que le terme général peut s'écrire sous la forme $u_{n}=u_{p}+(n-p) r$. On obtient :
\begin{enumMet}
    \item $u_{8}=u_{3}+(8-3) \times r \text { soit } 24=4+5 r \Longleftrightarrow 20=5 r \Longleftrightarrow r=4$ \\
    $u_{0}=u_{3}+(0-3) r \Longleftrightarrow u_{0}=4-3 \times 4=-8$ \\
    Donc $u_{n}=-8+4 n$
    \item $u_{9}=u_{5}+(9-5) \times r \text { soit } \dfrac{1}{4}=\dfrac{7}{4}+4 r \Longleftrightarrow-\dfrac{6}{4}=4 r \Longleftrightarrow r=-\dfrac{6}{16}=-\dfrac{3}{8}$\\
    $u_{0}=u_{5}+(0-5) r \Longleftrightarrow u_{0}=\dfrac{7}{4}-5 \times\left(-\dfrac{3}{8}\right)=\dfrac{29}{8}$ \\
    Donc $u_{n}=\dfrac{29}{8}-\dfrac{3}{8} n$ \\
    \item $u_{32}=u_{13}+(32-13) \times r_{\text {soit }}-7=16+19 r \Longleftrightarrow-23=19 r \Longleftrightarrow r=-\dfrac{23}{19}$\\
    $u_{0}=u_{13}+(0-13) r \Longleftrightarrow u_{0}=16-13 \times\left(-\dfrac{23}{19}\right)=\dfrac{603}{19}$\\
    Donc $u_{n}=\dfrac{603}{19}-\dfrac{23}{19} n$
    \item $u_{100}=u_{50}+(100-50) \times r$ soit $309=159+50 r \Longleftrightarrow 150=50 r \Longleftrightarrow r=3$.\\
    $u_{0}=u_{50}+(0-50) r \Longleftrightarrow u_{0}=159-50 \times 3=9$ \\
    Donc $u_{n}=9+3 n$

\end{enumMet}

\subsection*{Ex 60 p 34}

\begin{enumMet}
    \item D'après la formule $1+2+3+\cdots+n=\dfrac{n(n+1)}{2}$, on obtient :\\
    $S=\dfrac{73 \times 74}{2}=2701$
    \item $T=1+(1+3 \times 1)+(1+3 \times 2)+(1+3 \times 3)+\cdots+(1+3 \times 13)$\\
    $T=1 \times 14+3 \times(1+2+3+\cdots+13)$\\
    $T=14+3 \times \dfrac{13 \times 14}{2}=14+3 \times 91=287$\\
    \underline{\textbf{Autre méthode:}}\\
    Les termes $1 ; 4 ; 7 ; 10 ; \ldots$ sont les termes d'une suite arithmétique de premier terme $u_{0}=1$ et de raison $r=3$.\\
    La somme de $u_{0}+u_{1}+u_{2}+\cdots+u_{n}=(n+1) \dfrac{u_{0}+u_{n}}{2}$.\\
    On a donc la somme $T=u_{0}+u_{1}+u_{2}+u_{3}+\cdots+u_{13}$\\
    D'où $T=14 \times \dfrac{1+40}{2}=287$
    \item $U=(70+1)+(70+2)+(70+3)+\cdots+(70+30)$\\
    $U=70 \times 30+(1+2+3+\cdots+30)$\\
    $U=2100+\dfrac{30 \times 31}{2}=2100+465=2565$\\
    \underline{\textbf{Autre méthode:}}\\
    Les termes $71 ; 72 ; 73 ; \ldots$ sont les termes d'une suite arithmétique de premier terme 71 et de raison $1$.\\
    On a donc la somme $U=30 \times \dfrac{71+100}{2}=2565$
    \item $V=2 \times(1+2+3+\cdots+25)$\\
    $V=2 \times\left(\dfrac{25 \times 26}{2}\right)=650$\\
    \underline{\textbf{Autre méthode:}}\\
    Les termes $2 ; 4 ; 6 ; \ldots$ sont les termes d'une suite arithmétique de premier terme 2 et de raison 2 .\\
    On a donc la somme $V=25 \times \dfrac{2+50}{2}=650$
\end{enumMet}

\subsection*{Ex 62 p 35}
\begin{enumMet}
    \item $c_{1}=1,\quad c_{2}=2,\quad c_{3}=3,\quad c_{4}=4$\quad et\quad $c_{5}=5$.
    \item Pour tout entier $n \geqslant 1, c_{n+1}=c_{n}+1$. La suite est donc arithmétique de raison $r=1$ et de premier terme $c_{1}=1$
    \item On a donc, pour tout entier $n \geqslant 1, c_{n}=1+(n-1) \times 1=n$.
    \item Pour 7 niveaux, on a :\\
    $1+2+3+\cdots+7=\dfrac{7 \times 8}{2}=28$\\
    On utilise 28 canettes pour 7 niveaux.
    \item On résout $91=\dfrac{n(n+1)}{2}$\\
    $\Longleftrightarrow 182=n^{2}+n \Longleftrightarrow n^{2}+n-182=0 .$\\
    On calcule le discriminant : $\Delta=b^{2}-4 a c=1+4 \times 182=729$.\\
    Les solutions sont $n_{1}=\dfrac{-1-27}{2}=-14 \quad \text {et }\quad  n_{2}=\dfrac{-1+27}{2}=13$\\
    Seule la solution positive est valide. Il y aura donc 13 rangées.
\end{enumMet}

\subsection*{Ex 62 p 35}
\begin{enumMet}
    \item $c_{1}=1,\quad c_{2}=2,\quad c_{3}=3,\quad c_{4}=4$\quad et\quad $c_{5}=5$.
    \item Pour tout entier $n \geqslant 1, c_{n+1}=c_{n}+1$. La suite est donc arithmétique de raison $r=1$ et de premier terme $c_{1}=1$
    \item On a donc, pour tout entier $n \geqslant 1, c_{n}=1+(n-1) \times 1=n$.
    \item Pour 7 niveaux, on a :\\
    $1+2+3+\cdots+7=\dfrac{7 \times 8}{2}=28$\\
    On utilise 28 canettes pour 7 niveaux.
    \item On résout $91=\dfrac{n(n+1)}{2}$\\
    $\Longleftrightarrow 182=n^{2}+n \Longleftrightarrow n^{2}+n-182=0 .$\\
    On calcule le discriminant : $\Delta=b^{2}-4 a c=1+4 \times 182=729$.\\
    Les solutions sont $n_{1}=\dfrac{-1-27}{2}=-14 \quad \text {et }\quad  n_{2}=\dfrac{-1+27}{2}=13$\\
    Seule la solution positive est valide. Il y aura donc 13 rangées.
\end{enumMet}

\subsection*{Ex 61 p 35}
\begin{enumMet}
    \item 
    \begin{enumMet}
    \item $S=1+3+5+\cdots+31$\\
    $S=16 \times 1+2(1+2+\cdots+15)$
    $S=16+2 \times \dfrac{15 \times 16}{2}=256$
    \item $T=4+7+10+\cdots+19$\\
    $T=6 \times(-2)+3(2+3+\cdots+7)$\\
    $T=-12+3 \times\left(\dfrac{7 \times 8}{2}-1\right)=-12+81=69$
    \end{enumMet}
    \item 
    \begin{enumMet}
    \item $U=\sum \limits_{i=1}^{27} 3 i$\\
$U=3(1+2+3+\cdots+27)=3 \times \dfrac{27 \times 28}{2}=1134$
    \item $V=\sum \limits_{i=0}^{10}(5+4 i)$\\
$V=5 \times 11+4 \times \dfrac{10 \times 11}{2}=275$
    \end{enumMet}
\end{enumMet}

\subsection*{Ex 64 p 35}
\begin{enumMet}
    \item $s_{1}=500, s_{2}=600$ et $s_{3}=700$\\
    Pour tout entier $n \geqslant 1, s_{n+1}=s_{n}+100$.
    \item La suite est arithmétique de raison $r=100$ et de premier terme $u_{1}=500$ donc $u_{n}=500+(n-1) \times 100$ soit $u_{n}=100 n+400$.
    \item La somme totale épargnée est $S=s_{1}+s_{2}+\cdots+s_{n}=n \times \dfrac{s_{1}+s_{n}}{2} .$\\
    On doit résoudre $S \geq 4200$ $\Longleftrightarrow n \times \dfrac{500+100 n+400}{2} \geq 4200$\\
    $\Longleftrightarrow 100 n^{2}+900 n \geq 8400$\\
    $\Longleftrightarrow 100 n^{2}+900 n-8400 \geq 0$\\
    On calcule le discriminant
    $$
    \Delta=900^{2}+4 \times 8400=4170000
    $$
    Les racines sont :
    $$
    \begin{array}{l}
    n_{1}=\dfrac{-900-\sqrt{4170000}}{200}=\dfrac{-9-\sqrt{417}}{2}<0 \\
    n_{2}=\dfrac{-900+\sqrt{4170000}}{200}=\dfrac{-9+\sqrt{417}}{2} \approx 5,7
    \end{array}
    $$
    Le polynôme est du signe de $a=100$ donc positif à l'extérieur des racines.\\
    La famille pourra donc partir en voyage au bout de 6 ans.
\end{enumMet}

\section*{Suites géométriques}

\subsection*{Ex 29 p 31}
\begin{enumMet}
    \item 1. Pour tout entier, $n \in \mathbb{N}, v_{n+1}=4 \times v_{n}$ $v_{0}=3$\\
    $v_{1}=4 \times 3=12$\\
    $v_{2}=4 \times 12=48$\\
    $v_{3}=4 \times 48=192$\\
    $v_{4}=4 \times 192=768$
    \item Pour tout entier, $n \in \mathbb{N}, v_{n+1}=-3 \times v_{n}$ $v_{0}=2$\\
    $v_{1}=-3 \times 2=-6$\\
    $v_{2}=-3 \times(-6)=18$\\
    $v_{3}=-3 \times 18=-54$\\
    $v_{4}=-3 \times(-54)=162 v_{4}=-3 \times(-54)=162$
    \item Pour tout entier $n \in \mathbb{N}$, $v_{n+1}=\dfrac{1}{2} \times v_{n}$ $v_{0}=5$\\
    $v_{1}=\dfrac{1}{2} \times 5=\dfrac{5}{2}$\\
    $v_{2}=\dfrac{1}{2} \times \dfrac{5}{2}=\dfrac{5}{4}$\\
    $v_{3}=\dfrac{1}{2} \times \dfrac{5}{4}=\dfrac{5}{8}$\\
    $v_{4}=\dfrac{1}{2} \times \dfrac{5}{8}=\dfrac{5}{16}$
    \item Pour tout entier, $n \in \mathbb{N}, v_{n+1}=2 \times v_{n}$\\
    $v_{0}=-\dfrac{1}{4}$\\
    $v_{1}=2 \times\left(-\dfrac{1}{4}\right)=-\dfrac{1}{2}$\\
    $v_{2}=2 \times\left(-\dfrac{1}{2}\right)=-1$\\
    $v_{3}=2 \times(-1)=-2$\\
    $v_{4}=2 \times(-2)=-4$
\end{enumMet}

\subsection*{Ex 30 p 31}
\begin{enumMet}
    \item Pour tout entier $n \in \mathbb{N}, v_{n}=8^{n}$\\ 
    On calcule les premiers termes :
    $$
    \begin{array}{l}
    v_{0}=8^{0}=1 \\
    v_{1}=8^{1}=8 \\
    v_{2}=8^{2}=64
    \end{array}
    $$
    On a $\dfrac{v_{2}}{v_{1}}=\dfrac{64}{8}=8$ et $ \dfrac{v_{1}}{v_{0}}=\dfrac{8}{1}=8$. La suite semble donc géométrique. Il faut vérifier pour tout entier $n$.\\
    Pour tout entier $n \in \mathbb{N}, v_{n+1}=8^{n+1}=8^{n} \times 8=8 \times v_{n}$\\
    La suite est donc bien géométrique de raison $q=8$ et de premier terme $v_{0}=1$
    \item Pour tout entier $n \in \mathbb{N}, v_{n}=n^{4}$\\ 
    On calcule les premiers termes :\\
    $v_{0}=0^{4}=0$\\
    $v_{1}=1^{4}=1$\\
    $v_{2}=2^{4}=16$\\
    $v_{3}=3^{4}=81$\\
    On a $\dfrac{v_{3}}{v_{2}}=\dfrac{81}{16}$ et $\dfrac{v_{2}}{v_{1}}=\dfrac{16}{1}=16$\\
    La suite n'est donc pas géométrique.
    \item Pour tout entier $n \in \mathbb{N}, v_{n}=2 \times 3^{n}$\\
    On calcule les premiers termes :\\
    $v_{0}=2 \times 3^{0}=2$\\
    $v_{1}=2 \times 3^{1}=6$\\
    $v_{2}=2 \times 3^{2}=18$\\
    On a $\dfrac{v_{2}}{v_{1}}=\dfrac{18}{6}=3$ et $ \dfrac{v_{1}}{v_{0}}=\dfrac{6}{2}=3$. La suite semble donc géométrique. Il faut vérifier pour tout entier $n .$\\
    Pour tout entier $n \in \mathbb{N}, v_{n+1}=2 \times 3^{n+1}=2 \times 3^{n} \times 3=3 \times v_{n}$
    La suite est donc bien géométrique de raison $q=3$ et de premier terme $v_{0}=2$.
    \item Pour tout entier $n \in \mathbb{N}, v_{n}=-5 \times 2^{n+1}$\\
    On calcule les premiers termes : \\
    $v_{0}=-5 \times 2^{0+1}=-10$\\
    $v_{1}=-5 \times 2^{1+1}=-20$ \\
    $v_{2}=-5 \times 2^{2+1}=-40$ \\
    On a $\dfrac{v_{2}}{v_{1}}=\dfrac{-40}{-20}=2$ et $\dfrac{v_{1}}{v_{0}}=\dfrac{-20}{-10}=2$
    
    La suite semble donc géométrique.\\
    II faut vérifier pour tout entier $n$. Pour tout entier $n \in \mathbb{N},$\\
    $ v_{n+1}=-5 \times 2^{n+1+1}=-5 \times 2^{n+1} \times 2=2 \times v_{n}$\\
    La suite est donc bien géométrique de raison $q=2$ et de premier terme $v_{0}=-10$.
    \item $$
    \left\{\begin{array}{l}
    v_{0}=1 \\
    v_{n+1}=3 v_{n}+3 n \text {, pour tout } n \in \mathbb{N}
    \end{array}\right.
    $$
    On calcule les premiers termes :\\
    $v_{1}=3 v_{0}+3 \times 0=3 \times 1+3 \times 0=3$\\
    $v_{2}=3 v_{1}+3 \times 1=3 \times 3+3 \times 1=12$\\
    On a donc $\dfrac{v_{2}}{v_{1}}=\dfrac{12}{3}=4$ et $\dfrac{v_{1}}{v_{0}}=\dfrac{3}{1}=3$.\\
    La suite n'est donc pas géométrique.
    \item $$
    \left\{\begin{array}{l}
    v_{0}=2 \\
    v_{n+1}=5+7 v_{n}, \text { pour tout } n \in \mathbb{N}
    \end{array}\right.
    $$
    On calcule les premiers termes :
    $$
    \begin{array}{l}
    v_{1}=5+7 v_{0}=5+7 \times 2=19 \\
    v_{2}=5+7 v_{1}=5+7 \times 19=138 \\
    \text { On a donc } \overline{v_{1}}=\dfrac{138}{19} \approx 7,26 \dfrac{v_{1}}{\text { et } v_{0}}=\dfrac{19}{2} \approx 6,33
    \end{array}
    $$
    On a donc $\dfrac{v_{2}}{v_{1}}=\dfrac{138}{19} \approx 7,26$ et $\dfrac{v_{1}}{v_0}=\dfrac{19}{2} \approx 6,33$ La suite n'est donc pas géométrique.
\end{enumMet}

\subsection*{Ex 31 p 31}
\begin{enumMet}
    \item 
\end{enumMet}

\subsection*{Ex 69 p 36}
\begin{enumMet}
    \item On a $v_{n}=v_{0}\times q^{n}\iff v_3=v_0\times q^3 \iff 12=v_0\times 3^3 \iff v_0=\dfrac{12}{27}$\\
    Donc $v_{n}=\dfrac{12}{27}\times 3^{n}$ et $v_{20}=\dfrac{12}{27}\times3^{20}=1\ 549\ 681\ 956$
    \item On a $v_{n}=v_{0}\times q^{n}\iff v_{31}=v_0\times q^{31} \iff 32=v_0\times (-2)^{31} \iff v_0=-\dfrac{1}{67108864}$\\
    Donc $v_{n}=-\dfrac{1}{67108864}\times (-2)^{n}$ et $v_{20}=-\dfrac{1}{67108864}\times(-2)^{20}=-\dfrac{1}{64}$
    \item La suite $\left(v_{n}\right)$ est géométrique car $v_{n+1}=2 v_{n}$. La raison est donc $q=2$ et le premier terme $v_{0}=-5$\\
    On a $v_{n}=v_{0}\times q^{n}\iff v_{n}=-5\times 2^{n}$\\
    Donc $v_{20}=-5\times2^{20}=-5\ 242\ 880$
    \item La suite $\left(v_{n}\right)$ est géométrique car $v_{n+1}=-\dfrac{1}{2} v_{n}$. La raison est donc $q=-\dfrac{1}{2}$ et le terme $v_{1}=2048$, donc $v_0=-4096$.\\
    On a $v_{n}=v_{0}\times q^{n}\iff v_{n}=-4096\times \left(-\dfrac{1}{2}\right)^{n}$\\
    Donc $v_{20}=-4096\times \left(-\dfrac{1}{2}\right)^{20}=-\dfrac{1}{256}$
\end{enumMet}

\subsection*{Ex 70 p 36}
\begin{enumMet}
    \item $v_{n}=\dfrac{4^{n}}{3^{n+1}}$\\
    On calcule les premiers termes:\\
    $v_{0}=\dfrac{4^{0}}{3^{0+1}}=\dfrac{1}{3}$\\
    $v_{1}=\dfrac{4^{1}}{3^{1+1}}=\dfrac{4}{9}$\\
    $v_{2}=\dfrac{4^{2}}{3^{2+1}}=\dfrac{16}{27}$\\
    On a $\dfrac{v_{2}}{v_{1}}=\dfrac{\dfrac{16}{27}}{\dfrac{4}{9}}=\dfrac{4}{3}$ et $\dfrac{v_{1}}{v_{0}}=\dfrac{\dfrac{4}{9}}{\dfrac{1}{3}}=\dfrac{4}{3}$. La suite semble donc géométrique.\\
    Il faut vérifier : Pour tout entier $n$,\\
    $v_{n+1}=\dfrac{4^{n+1}}{3^{n+1+1}}=\dfrac{4^{n} \times 4}{3^{n+1} \times 3}=\dfrac{4}{3} v_{n}$\\
    La suite est donc bien géométrique de raison $q=\dfrac{4}{3}$ et de premier terme $v_{0}=\dfrac{1}{3}$
    \item $v_{n}=(-7)^{n}$\\
    On calcule les premiers termes:\\
    $v_{0}=(-7)^{0}=1$\\
    $v_{1}=(-7)^{1}=-7$\\
    $v_{2}=(-7)^{2}=49$\\
    On a $\dfrac{v_{2}}{v_{1}}=\dfrac{-49}{7}=-7$ et $\dfrac{v_{1}}{v_{0}}=\dfrac{-7}{1}=-7$.\\
    La suite semble donc géométrique. Il faut vérifier : Pour tout entier $n$,\\
    $v_{n+1}=(-7)^{n+1}=(-7)^{n} \times(-7)=-7 v_{n}$
    La suite est donc bien géométrique de raison $q=-7$ et de premier terme $v_{0}=1$
    \item $v_{n}=5 n+2^{n}$\\
    On calcule les premiers termes :
    $$
    \begin{array}{l}
    v_{0}=5 \times 0+2^{0}=1 \\
    v_{1}=5 \times 1+2^{1}=7 \\
    v_{2}=5 \times 2+2^{2}=14
    \end{array}
    $$
    On $\dfrac{v_{2}}{v_{1}}=\dfrac{14}{7}=2 \dfrac{v_{1}}{v_0}=\dfrac{7}{1}=7$
    La suite n'est donc pas géométrique.
    \item $v_{n}=\dfrac{1}{3^{n}}$\\
    On calcule les premiers termes :
    $$
    \begin{array}{l}
    v_{0}=\dfrac{1}{3^{0}}=1 \\
    v_{1}=\dfrac{1}{3^{1}}=\dfrac{1}{3} \\
    v_{2}=\dfrac{1}{3^{2}}=\dfrac{1}{9}
    \end{array}
    $$
    On a $\dfrac{v_{2}}{v_{1}}=\dfrac{\dfrac{1}{9}}{\dfrac{1}{3}}=\dfrac{1}{3} \quad$ et $\dfrac{v_{1}}{v_{0}}=\dfrac{\dfrac{1}{3}}{1}=\dfrac{1}{3}$.\\
    La suite semble donc géométrique. II faut vérifier : Pour tout entier $n$,\\
    $v_{n+1}=\dfrac{1}{3^{n+1}}=\dfrac{1}{3^{n}} \times \dfrac{1}{3}=\dfrac{1}{3} v_{n}$\\
    La suite est donc bien géométrique de raison $q=\dfrac{1}{3}$ et de premier terme $v_{0}=1$.
\end{enumMet}

\subsection*{Ex 71 p 36}
La suite étant géométrique, on sait que le terme général peut s'écrire sous la forme $v_{n}=v_{p} \times q^{n-p}$ On obtient :
\begin{enumMet}
    \item $v_{8}=v_{3} \times q^{8-3} \text { soit } 1458=6 \times q^{5} \iff q^{5}=243 \iff q=3$.\\
    $v_{0}=v_{3} \div q^{3}=\dfrac{6}{3^{3}}=\dfrac{2}{9}$\\
    Donc, pour tout entier naturel $n$, $v_{n}=\dfrac{2}{9} \times 3^{n}$
    \item $v_{12}=v_{6} \times q^{12-6} \text { soit }-\dfrac{9}{32}=-18 \times q^{6} \iff q^{6}=\dfrac{-\dfrac{9}{32}}{-18}=\dfrac{1}{64} \iff q=\dfrac{1}{2} \text { ou } q=-\dfrac{1}{2}$\\
    \begin{itemMet}
        \item Si $q=\dfrac{1}{2}$\\
        $v_{0}=v_{6} \div q^{6}=-18 \div\left(\dfrac{1}{2}\right)^{6}=-18 \times 2^{6}=-1152$\\
        Donc, pour tout entier naturel $n,$
        $v_{n}=-1152 \times\left(\dfrac{1}{2}\right)^{n}$\\
        \item Si $q=-\dfrac{1}{2}$\\
        $v_{0}=v_{6} \div q^{6}=-18 \div\left(-\dfrac{1}{2}\right)^{6}=-1152$\\
        Donc, pour tout entier naturel $n$, $v_n=1152 \times\left(-\dfrac{1}{2}\right)^{n}$
    \end{itemMet}
    \item $v_{15}=v_{10} \times q^{15-10}$ soit $46875=15 \times q^{5} \iff q^{5}=3125 \iff q=5$.\\
    $v_{0}=v_{10} \div q^{10}=\dfrac{15}{5^{10}}=\dfrac{3}{5^{9}}$\\
    Donc, pour tout entier naturel $n$, $v_{n}=\dfrac{3}{5^{9}} \times 5^{n}=3 \times 5^{n-9}$
    \item $v_{23}=v_{21} \times q^{23-21}$ soit $262144=65536 \times q^{2} \iff q^{2}=4 \iff q=2$ ou $q=-2$.
    \begin{itemMet}
        \item Si $q=2, v_{0}=v_{21} \div q^{21}=65536 \div 2^{21}=\dfrac{1}{32}$\\
        Donc, pour tout entier naturel $n$ $v_{n}=\dfrac{1}{32} \times 2^{n}$
        \item Si $q=-2$,$v_{0}=v_{21} \div q^{21}=65536$\\
        Donc, pour tout entier naturel $n$, $v_{n}=-\dfrac{1}{32} \times(-2)^{n}$
    \end{itemMet}
\end{enumMet}

\subsection*{Ex 72 p 36}
D'après la formule $1+q+q^{2}+\cdots+q^{n}=\dfrac{1-q^{n+1}}{1-q}$, on obtient :
\begin{enumMet}
    \item $S=1+4+4^{2}+\cdots+4^{9}=\dfrac{1-4^{10}}{1-4}$\\
    $S=\dfrac{1048576-1}{3}=349\ 525$
    \item $T=3(1-2+4-8+\cdots+64)$\\
    $T=3\left(1+(-2)^{1}+(-2)^{2}+(-2)^{3}+\cdots+(-2)^{6}\right)$\\
    $T=3 \times \dfrac{1-(-2)^{7}}{1-(-2)}$\\
    $T=3 \times \dfrac{129}{3}=129$
    \item $U=9\left(1+\dfrac{1}{3}+\dfrac{1}{3^{2}}+\dfrac{1}{3^{3}} \cdots+\dfrac{1}{3^{8}}\right)$\\
    $U=9 \times \dfrac{1-\left(\dfrac{1}{3}\right)^{9}}{1-\dfrac{1}{3}}=\dfrac{9841}{729}$
    \item $V=1+\dfrac{1}{2}+\dfrac{1}{2^{2}}+\cdots+\dfrac{1}{2^{5}}=\dfrac{1-\left(\dfrac{1}{2}\right)^{6}}{1-\dfrac{1}{2}}=\dfrac{63}{32}$
\end{enumMet}

\subsection*{Ex 73 p 36}
\begin{enumMet}
    \item 
    \begin{enumMet}
        \item $S=1+2+2^{2}+2^{3}+\cdots+2^{15}=\dfrac{1-2^{16}}{1-2}=65535$
        \item $T=\left(\dfrac{1}{4}\right)^{2}+\left(\dfrac{1}{4}\right)^{3}+\cdots+\left(\dfrac{1}{4}\right)^{7}$\\
        $T=\left(\dfrac{1}{4}\right)^{2}\left(1+\dfrac{1}{4}+\cdots+\dfrac{1}{4^{5}}\right)=\dfrac{1}{16} \times \dfrac{1-\dfrac{1}{4^{6}}}{1-\dfrac{1}{4}}$\\
        $T=\dfrac{1365}{16384}$
    \end{enumMet}
    \item
    \begin{enumMet}
        \item $U=\sum \limits_{i=0}^{8}\left(-\dfrac{1}{2}\right)^{i}$\\
        $U=\dfrac{1-\left(-\dfrac{1}{2}\right)^{9}}{1+\dfrac{1}{2}}$\\
        $U =\dfrac{171}{256}$
        \item $V=\sum_{i=0}^{8}\left(2 \times 3^{i}\right)$\\
        $V=2\left(1+3+9+27+\cdots+3^{8}\right)=2 \times \dfrac{1-3^{9}}{1-3}=19\ 682$
    \end{enumMet}
\end{enumMet}

\subsection*{Ex 74 p 36}
\begin{enumMet}
    \item Le deuxième jour, on a :\\
$20+\dfrac{15}{100} \times 20=20 \times\left(1+\dfrac{15}{100}\right)=20 \times 1,15=23$ mm$^2$ de peau.\\
Le troisième jour, on a :\\
$23+\dfrac{15}{100} \times 23=23 \times 1,15=26,45$ mm$^2$ de peau.\\
\item Pour tout $n \in \mathbb{N}, v_{n+1}=v_{n}+\dfrac{15}{100} v_{n}=1,15 v_{n}$.
\item La suite est donc géométrique de raison $q=1,15$ et de premier terme $v_{0}=20$.
\item On a donc, pour tout $n \in \mathbb{N}, v_{n}=20 \times 1,15^{n}$.
\item On doit résoudre $v_{n} \geq 1000 \iff 20 \times 1,15^{n} \geq 1000 \iff 1,15^{n} \geq 50$.\\
Avec la calculatrice, on trouve $1,15^{27} \approx 43,54$ et $1,15^{28}=50,07$.\\
Il faut donc 28 jours pour pouvoir greffer la surface nécessaire.
\end{enumMet}

\subsection*{Ex 80 p 37}
%\begin{enumMet}
    $v_{n}$ modélise la distance parcourue au jour $n$. Le premier terme est $v_{1}=20$ et, pour tout entier $n \geqslant 1,$ on a $v_{n+1}=0,95 v_{n}$\\
    La suite est géométrique de premier terme $v_{1}=20$ et de raison $q=0,95$.\\
    On cherche $n$ tel que $v_{1}+v_{2}+\cdots+v_{n}>200$.\\
    $v_{1}+v_{2}+\ldots v_{n}=v_{1} \times \dfrac{1-0,95^{n}}{1-0,95}=20 \times \dfrac{1-0,95^{n}}{0,05}=400\left(1-0,95^{n}\right)$\\
    $400\left(1-0,95^{n}\right)>200 \Longleftrightarrow 1-0,95^{n}>0,5$\\
    $\Longleftrightarrow 0,95^{n}<0,5$\\
    À l'aide de la calculatrice, on obtient $n=14$ La randonnée sera terminée en 14 jours.
%\end{enumMet}

\section*{variations de suites}

\subsection*{21 p 31}
La suite $\left(u_{n}\right)$ est géométrique de raison $\left.\dfrac{2}{3} \in\right] 0 ; 1[$ donc elle est décroissante sur $\mathbb{N}$.\\
La suite $\left(v_{n}\right)$ est géométrique de raison $\dfrac{7}{3}>1$ donc elle est croissante sur $\mathbb{N}$.

\subsection*{24 p 31}
\begin{enumMet}
    \item \ \\
    $
    \begin{aligned}
    u_{0} &=20 \\
    u_{1} &=u_{0}-(0-4)^{2}=20-16=4 \\
    u_{2} &=u_{1}-(1-4)^{2}=4-9=-5 \\
    u_{3} &=u_{2}-(2-4)^{2}=-5-4=-9
    \end{aligned}
    $
    \item La suite semble être décroissante.
    \item Pour tout, $n \in \mathbb{N}, u_{n+1}-u_{n}=-(n-4)^{2}<0$ donc $u_{n+1}<u_{n}$ ce qui montre bien que la suite $\left(u_{n}\right)$ est décroissante.
\end{enumMet}

\subsection*{38 p 32}
Pour tout entier $n \in \mathbb{N}, u_{n+1}=\dfrac{1}{4(n+1)+3}=\dfrac{1}{4 n+7}$.\\
Pour tout entier $n \in \mathbb{N}$, $u_{n+1}-u_{n}=\dfrac{1}{4 n+7}-\dfrac{1}{4 n+3}$ soit\\
$u_{n+1}-u_{n}=\dfrac{4 n+3}{(4 n+7)(4 n+3)}-\dfrac{4 n+7}{(4 n+3)(4 n+7)}$\\
$u_{n+1}-u_{n}=\dfrac{-4}{(4 n+7)(4 n+3)}$\\
Comme $n \geqslant 0$, $4 n+3 \geqslant 0$ et $4 n+7 \geqslant 0$

\subsection*{39 p 32}
\begin{enumMet}
    \item La suite $\left(u_{n}\right)$ est géométrique de raison $q=2,5$ supérieure à 1 et de premier terme 1 qui est positif, donc elle est croissante.
    \item La suite $\left(v_{n}\right)$ est géométrique de raison $q=0,4$ comprise entre 0 et 1 et de premier terme 1 qui est positif, donc elle est décroissante.
    \item La suite $\left(w_{n}\right)$ est géométrique de raison $q=-1,8$ strictement inférieure à 0 et de premier terme 1 qui est positif, donc elle n'est pas monotone.
    \item Pour tout entier $n \in \mathbb{N}, t_{n+1}=3(n+1)+5=3 n+8$
    Pour tout entier $n \in \mathbb{N}, t_{n+1}-t_{n}=3 n+8-(3 n+5)=3$
    La suite $\left(t_{n}\right)$ est arithmétique de raison $r=3$ qui est supérieure à $0,$ donc elle est croissante.\\
    \grasInMet{Remarque :} On peut également reconnaître que $u_{n}$ est de la forme $u_{0}+n r$ et déduire directement que $\left(u_{n}\right)$ est arithmétique de raison $r=3$ et de premier terme $u_{0}=5$.
\end{enumMet}

\subsection*{49 p 31}
\begin{enumMet}
    \item Pour tout entier naturel $n$ :\\
    $$u_{n+1}=\dfrac{n+1+1}{n+1+2}=\dfrac{n+2}{n+3}$$
    $$
    u_{n+1}-u_{n}=\dfrac{n+2}{n+3}-\dfrac{n+1}{n+2}=\dfrac{(n+2)^{2}-(n+1)(n+3)}{(n+2)(n+3)}
    $$
    $$u_{n+1}-u_{n}=\dfrac{n^{2}+4 n+4-\left(n^{2}+4 n+3\right)}{(n+2)(n+3)}$$
    $$u_{n+1}-u_{n}=\dfrac{1}{(n+2)(n+3)}$$
    Comme $n \geq 0$ alors $n+2 \geq 2$ et $n+3 \geq 3, u_{n+1}-u_{n}>0$ La suite est donc croissante.
    \item On a $u_{n}>0$ pour tout entier naturel, $n \geqslant 1$, $$u_{n+1}=\dfrac{3^{n+1}}{n+1}$$
    $$
    u_{n+1}-u_{n}=\dfrac{3^{n+1}}{n+1}-\dfrac{3^{n}}{n}=\dfrac{3^{n+1} \times n}{n(n+1)}-\dfrac{3^{n}(n+1)}{n(n+1)}
    $$
    $$u_{n+1}-u_{n}=\dfrac{3^{n} \times 3 \times n-3^{n}(n+1)}{n(n+1)}=\dfrac{3^{n}(3 n-(n+1))}{n(n+1)}=\dfrac{\left.3^{n}(3 n-n-1\right)}{n(n+1)}=$$
    $$\dfrac{3^{n}(2 n-1)}{n(n+1)}$$
    Or $3^{n}>0$ et comme $n \geq 1, n(n+1)>0$. Donc $u_{n+1}-u_{n}$ est du signe de $2 n-1$. Comme $n>1$, on a donc $2 n>2$ et donc $2 n-1>1>0$. D'où $u_{n+1}-u_{n}>0$. La suite est donc croissante.
    \item Pour tout entier naturel $n$ :
    $$
    u_{n+1}=(n+1)^{2}-3(n+1)+12=n^{2}+2 n+1-3 n-3+12=n^{2}-n+10
    $$
    $$u_{n+1}-u_{n}=n^{2}-n+10-\left(n^{2}-3 n+12\right)=2 n-2$$
    On a $2 n-2>0$ si $2 n>2 \Longleftrightarrow n>1$. La suite est donc croissante à partir de $n=1$.
    \item Pour tout entier naturel $n \geqslant 1$ :
    $$
    u_{n}=\dfrac{n+1}{n(n+1)}-\dfrac{n}{n(n+1)}=\dfrac{1}{n(n+1)}
    $$
    Les termes de la suite sont donc strictement positifs. $$u_{n+1}=\dfrac{1}{(n+1)(n+2)}$$
    $$
    \text { D'où } \dfrac{u_{n+1}}{u_{n}}=\dfrac{1}{(n+1)(n+2)} \times n(n+1)=\dfrac{n}{n+2}
    $$
    Or $n+2>n$ donc $\dfrac{n}{n+2}<1$\\
    Par conséquent, pour tout entier $n \geqslant 1,$ on a $u_{n}>0$ et $n \geqslant 1,$ on a $u_{n}>0$ et $\dfrac{u_{n+1}}{u_{n}}<1$ La suite est donc décroissante
\end{enumMet}

\subsection*{65 p 35}
\begin{enumMet}
    \item \begin{enumMet}
        \item La suite arithmétique de raison $4>0$ donc la suite est croissante.
        \item La suite est arithmétique de raison $-\dfrac{3}{8}<0$ donc la suite est décroissante.
        \item $u_{n}=\dfrac{n^{2}+4 n+3}{n+3}=\dfrac{(n+1)(n+3)}{n+3}=n+1$
        La suite est arithmétique de raison $1>0$ donc la suite est croissante.
        \item $u_{n}=\dfrac{3 n^{2}+5 n-2}{n+2}=\dfrac{(n+2)(3 n-1)}{n+2}=3 n-1$
        La suite est arithmétique de raison $3>0$ donc la suite est croissante.
    \end{enumMet}
    \item La suite étant arithmétique, on sait que le terme général peut s'écrire sous la forme $u_{n}=u_{p}+(n-p) r$. On obtient
    \begin{enumMet}
        \item $u_{8}=u_{3}+(8-3) \times r$ soit $24=4+5 r \Longleftrightarrow 20=5 r \Longleftrightarrow r=4$\\
        La raison étant $r=4>0,$ la suite est croissante.\\
        \item $u_{9}=u_{5}+(9-5) \times r$ soit
        $\dfrac{1}{4}=\dfrac{7}{4}+4 r \Longleftrightarrow-\dfrac{6}{4}=4 r \Longleftrightarrow r=-\dfrac{6}{16}=-\dfrac{3}{8}$\\
        La raison étant $r=-\dfrac{3}{8}<0$, la suite est décroissante.
        \item $u_{32}=u_{13}+(32-13) \times r_{\text {soit }}$
        $$
        -7=16+19 r \Longleftrightarrow-23=19 r \Longleftrightarrow r=-\dfrac{23}{19}
        $$
        La raison étant $r=-\dfrac{23}{19}<0$, la suite est décroissante.
        \item $u_{100}=u_{50}+(100-50) \times r$ soit\\
        $309=159+50 r \Longleftrightarrow 150=50 r \Longleftrightarrow r=3$\\
        La raison étant $r=3,$ la suite est croissante.
    \end{enumMet}
\end{enumMet}

\subsection*{77 p 37}
\begin{enumMet}
    \item \begin{enumMet}
        \item $v_{n}=\dfrac{4^{n}}{3^{n+1}}$\\
        Pour tout entier $n$ :\\
        $v_{n+1}=\dfrac{4^{n+1}}{3^{n+1+1}}=\dfrac{4^{n} \times 4}{3^{n+1} \times 3}=\dfrac{4}{3} v_{n}$\\
        La suite est géométrique de raison $q=\dfrac{4}{3}$ et de premier terme $v_{0}=\dfrac{1}{3}$\\
        Comme $q>1$ et $v_{0}>0$, la suite est croissante.
        \item $v_{n}=(-7)^{n}$\\
        Pour tout entier $n$ :\\
        $v_{n+1}=(-7)^{n+1}=(-7)^{n} \times(-7)=-7 v_{n}$\\
        La suite est géométrique de raison $q=-7$ et de premier terme $v_{0}=1$\\
        Comme $q<0$, la suite n'est pas monotone.\\
        \item $v_{n}=5^{n} \times 2^{n}$\\
        Donc $v_{n}=10^{n},$ qui est une suite géométrique de raison $q=10$ et de premier terme $v_{0}=1$\\
        Comme $q>1$, et $v_{0}>0$, la suite est croissante.
        \item $v_{n}=\dfrac{1}{3^{n}}$\\
        Pour tout entier $n$ :\\
        $v_{n+1}=\dfrac{1}{3^{n+1}}=\dfrac{1}{3^{n}} \times \dfrac{1}{3}=\dfrac{1}{3} v_{n}$\\
        La suite est géométrique de raison $q=\dfrac{1}{3}$ et de premier terme $v_{0}=1$.\\
        Comme $0<q<1$, et $v_{0}>0$, la suite est décroissante.
    \end{enumMet}
    \item La suite étant géométrique, on sait que le terme général peut s'écrire sous la forme $v_{n}=v_{p} \times q^{n-p}$. On obtient :
    \begin{enumMet}
        \item $v_{8}=v_{3} \times q^{8-3}$ soit $1458=6 \times q^{5} \Longleftrightarrow q^{5}=243 \Longleftrightarrow q=3$\\
        $v_{0}=v_{3} \times q^{0-3}=\dfrac{6}{3^{3}}=\dfrac{2}{9}$\\
        Comme $q>1$, et $v_{0}>0$, la suite est croissante.
        \item $v_{12}=v_{6} \times q^{12-6}$ soit $-\dfrac{9}{32}=-18 \times q^{6} \Longleftrightarrow q^{6}=\dfrac{-\dfrac{9}{32}}{-18}=\dfrac{1}{64} \Longleftrightarrow q=\dfrac{1}{2}$ ou $q=-\dfrac{1}{2}$.\\
        $v_{0}=v_{6} \times q^{0-6}=-18 \times\left(\dfrac{1}{2}\right)^{-6}=-18 \times 2^{6}=-1152$
        \begin{itemMet}
            \item si $q=\dfrac{1}{2}$\\
            Comme $0<q<1$, et $v_{0}<0$, la suite est croissante
            \item si $q=-\dfrac{1}{2}$\\
            Comme $q<0$, la suite n'est pas monotone.
        \end{itemMet}
        \item $v_{15}=v_{10} \times q^{15-10}$ soit $46875=15 \times q^{5} \Longleftrightarrow q^{5}=3125 \Longleftrightarrow q=5$\\
        $v_{0}=v_{10} \times q^{0-10}=\dfrac{15}{5^{10}}=\dfrac{3}{5^{9}}$\\
        Comme $q>1$, et $v_{0}>0,$ la suite est croissante.
        \item $v_{23}=v_{21} \times q^{23-21}$ soit $262144=65536 \times q^{2} \Longleftrightarrow q^{2}=4 \Longleftrightarrow q=2$ ou $q=-2$\\
        \begin{itemMet}
            \item Si $q=2, v_{0}=v_{21} \times q^{0-21}=65536 \times 2^{-21}=\dfrac{1}{32}$\\
            Comme $q>2$, et $v_{0}>0$, la suite est croissante.
            \item Si $q=-2, v_{0}=v_{21} \times q^{0-21}=65536 \times(-2)^{-21}=-\dfrac{1}{32}$\\
            Comme $q<0,$ la suite n'est pas monotone.
        \end{itemMet}
         
    \end{enumMet}
\end{enumMet}

\section*{Synthèse}

\subsection*{Ex 34 p 32}
\begin{enumMet}
    \item On calcule les premiers termes:\\
    $u_{0}=5, u_{1}=7, u_{2}=9$\\
    Comme $u_{1}-u_{0}=2$ et $u_{2}-u_{1}=2,$ la suite semble être arithmétique. On vérifie pour tout entier naturel $n$ :\\
    $u_{n+1}-u_{n}=2(n+1)+5-(2 n+5)$\\
    $u_{n+1}-u_{n}=2 n+2+5-2 n-5=2$\\
    La suite $\left(u_{n}\right)$ est donc arithmétique de raison $r=2$ et de premier terme $u_{0}=5$.\\
    \grasInMet{Remarque :} On peut également reconnaître que $u_{n}$ est de la forme $u_{0}+n r$ et déduire directement que $\left(u_{n}\right)$ est arithmétique de raison $r=2$ et de premier terme $u_{0}=5$.
    \item On calcule les premiers termes:
    $v_{0}=1$ donc $v_{1}=3 \times 1+2=5$ et $v_{2}=3 \times 5+2=17$
    $v_{1}-v_{0}=4$ et $v_{2}-v_{1}=12,$ la suite n'est donc pas arithmétique. $\dfrac{v_{1}}{v_{0}}=5 \dfrac{v_{2}}{\text { et } v_{1}}=\dfrac{17}{5}$, la suite n'est donc pas géométrique.
    \item 3. On calcule les premiers termes:
    $$
    w_{0}=\left(\dfrac{2}{3}\right)^{0}=1 \quad w_{1}=\left(\dfrac{2}{3}\right)^{1}=\dfrac{2}{3} w_{2}=\left(\dfrac{2}{3}\right)^{2}=\dfrac{4}{9}
    $$
    Comme $\dfrac{w_{1}}{w_{0}}=\dfrac{2}{3}$ et $\dfrac{w_{2}}{w_{1}}=\dfrac{2}{3},$ la suite semble géométrique. On vérifie pour tout entier naturel $n$ :
    $$
    \dfrac{w_{n+1}}{w_{n}}=\dfrac{\left(\dfrac{2}{3}\right)^{n+1}}{\left(\dfrac{2}{3}\right)^{n}}=\dfrac{2}{3}
    $$
    La suite $\left(w_{n}\right)$ est donc géométrique de raison $q=\dfrac{2}{3}$ et de premier terme $w_{0}=1$.\\
    \grasInMet{Remarque :} On peut également reconnaître la forme $q^{n}$ et déduire que la suite est géométrique de raison $q=\dfrac{2}{3}$ et de premier terme $w_{0}=1$.
    \item Comme on a $t_{n+1}=7 t_{n}$ pour tout entier $n \in \mathbb{N}$, la suite est géométrique de raison $q=7$ et de premier terme $t_{0}=1$
    \item Comme on a $s_{n+1}=s_{n}+4$ pour tout entier $n \in \mathbb{N},$ la suite est arithmétique de raison $n \in \mathbb{N},$ la suite est arithmétique de raison $r=4$ et de premier terme $s_{0}=1 .$
\end{enumMet}

\subsection*{Ex 37 p 32}
\begin{enumMet}
    \item $S=2(1+2+3+4+5+\cdots+25)$ \\
    $S=2 \times \dfrac{25 \times 26}{2}=650$ en utilisant la formule $1+2+\ldots+n=\dfrac{n(n+1)}{2}$
    \item $T=1+2^{1}+2^{2}+2^{3}+\cdots+2^{13}$\\
    $T=\dfrac{1-2^{14}}{1-2}=2^{14}-1=16383$ en utilisant la formule
    $1+q+q^{2}+\ldots+q^{n}=\dfrac{1-q^{n+1}}{1-q}$
\end{enumMet}

\subsection*{Ex 83 p 38}
\begin{enumMet}
    \item 
\end{enumMet}

\subsection*{Ex 84 p 38}
\begin{enumMet}
    \item 
\end{enumMet}

\subsection*{Ex 85 p 38}
\begin{enumMet}
    \item 
\end{enumMet}


\end{document}

